\documentclass[a4paper,12pt]{report}

%  PREAMBOLO 
\usepackage[a4paper]{geometry}
\usepackage{amssymb,amsmath,amsthm}
\usepackage{graphicx}
\usepackage{url}
\usepackage{epsfig}
\usepackage[italian]{babel}
\usepackage{setspace}
\usepackage{tesi} 
\usepackage{wrapfig}
\usepackage[hidelinks]{hyperref}
\usepackage[export]{adjustbox}
\usepackage{epigraph}
\usepackage{blindtext}
\usepackage{subfig} % per le figure con più pannelli (a, b, c)
\usepackage{booktabs} % per le tabelle formattate correttamente (senza linee verticali) 
\usepackage{tikz}
\usetikzlibrary{positioning, arrows.meta, fit, backgrounds, calc}
\input{tikz_common_styles}   


\usepackage[utf8]{inputenc}
\setlength{\marginparwidth}{2cm}
\usepackage[disable]{todonotes} %per note e commenti

\usepackage{csquotes}
%  BIBLIOGRAFIA: motore biblatex/biber per citazioni a piè di pagina 
\usepackage[backend=biber, citestyle=verbose-ibid, bibstyle=numeric, sorting=none, giveninits=true]{biblatex}
\addbibresource{bibliografia.bib}

\newtheorem{myteor}{Teorema}[section]
\newenvironment{teor}{\begin{myteor}\sl}{\end{myteor}}

%  NUMERAZIONE CAPITOLI IN LETTERE (richiesta Prof. Buffa: Primo, secondo ecc.
\usepackage{titlesec}

% Creiamo un convertitore invisibile per il titolo
\newcommand{\capitoloInLettere}{%
  \ifcase\value{chapter}\or Primo\or Secondo\or Terzo\or Quarto\or Quinto\or Sesto\or Settimo\or Ottavo\or Nono\or Decimo\fi
}

% Ridisegniamo la grafica del titolo della pagina
\titleformat{\chapter}[display]
  {\normalfont\huge\bfseries}
  {\chaptertitlename\ \capitoloInLettere} % Qui stampa la magia: "\hyperref[cap1]{Capitolo Primo}"
  {20pt}
  {\Huge}

%  TITOLO E DATI 
\begin{document}
\title{Sistemi d'Arma Autonomi Letali: \\criticità nel trattamento dei dati sensibili e analisi sperimentale tra vulnerabilità algoritmiche e giuridiche}
\author{Daniele Intra}
\dept{Corso di Laurea Magistrale in Sicurezza Informatica} 
\anno{2025-2026}
\matricola{53665A}
\relatore{Prof. Matteo Buffa}
\correlatore{Prof. Massimo Walter Rivolta}

%  DEDICA E PREFAZIONE 
\beforepreface

\prefacesection{Dedica}
\begin{flushright}
\begin{minipage}{1\textwidth}
\Large\sl
Dedicato alle persone che amo, a tutti quelli che mi hanno affiancato in questo momento difficile della mia vita, alla mia famiglia, a mio nonno e mia nonna.

Dedicato a Manuel Pecchenini, che mi ha insegnato la forza di non arrendersi mai, a non mollare mai e a lottare per ciò in cui si crede.

Dedicato a chi lotta per un mondo più giusto e umano, ai milioni di bambini che soffrono ogni secondo in tutto il mondo, a chi non ha voce e a chi non ha diritto di parola.

Dedicato a tutte le studentesse e agli studenti che non sono più qui con noi e che non hanno potuto scrivere una loro dedica, dedicato alla loro lotta contro un sistema scolastico e universitario che non li ha mai ascoltati e che li ha oppressi, rendendoli un corto ed invisibile servizio da telegiornale.

Dedicato ai miei studenti, ex studenti, colleghi e alla Scuola, che mi ha permesso di crescere e di diventare la persona che sono oggi.

Dedicato a chi crede che la tecnologia possa essere un mezzo per il progresso e non una mera sostituzione di tutto quello che ci rende umani.
\end{minipage}
\end{flushright}

\prefacesection{Prefazione}
\begin{flushright}
\begin{minipage}{0.6\textwidth}
\epigraph{
    ``Il governo dei manganelli e dei plotoni di esecuzione, della carestia artificiale, dell'imprigionamento in massa e della deportazione di massa, non solo è inumano (nessuno se ne preoccupa più di tanto ai giorni nostri), ma è palesemente inefficiente, e in un'epoca di tecnologia avanzata l'inefficienza è un peccato mortale. Uno Stato totalitario davvero efficiente sarebbe quello in cui l'onnipotente potere esecutivo dei capi politici e il loro corpo manageriale controllano una popolazione di schiavi che non devono essere costretti ad esserlo con la forza perché amano la loro schiavitù.''
}{\textit{Aldous Huxley, Brave New World (1932)}}
\end{minipage}
\end{flushright}

Questo lavoro nasce da una profonda necessità di comprendere e resistere all'irrefrenabile insinuazione della \textit{technè} nella sfera più intima e vulnerabile dell'essere umano. In una società che chiede insistentemente di essere efficienti, produttivi, iperconnessi e miscelati con le macchine, occorre interrogarsi sulle implicazioni sociali e politiche di questa trasformazione.
L'intelligenza artificiale in ambito militare ha infatti cessato di essere un semplice \textit{strumento} per diventare un vero e proprio \textit{agente}, spingendo il dibattito pubblico a concentrarsi quasi esclusivamente su dilemmi filosofici. Ma l'etica da sola si rivela inefficace se non viene tradotta in rigidi vincoli ingegneristici.
Rispetto a una canonica tesi in sicurezza informatica, questo lavoro si focalizza sulla dimensione umanitaria e giuridica per dimostrare una tesi ben precisa: la violazione dei diritti umani, nel contesto dei Lethal Autonomous Weapon Systems (LAWS), non deriva solo da ciniche scelte geopolitiche, ma anche e soprattutto da vulnerabilità architetturali intrinseche ai sistemi stessi. Criticità come i bias nei dati di addestramento, la non interpretabilità delle reti neurali (\textit{Black Box}), la suscettibilità al \textit{Data Poisoning} e agli \textit{Adversarial Attacks} si trasformano direttamente in danni collaterali sul campo.
Attraverso un'analisi interdisciplinare, la tesi si propone quindi di esaminare il trattamento dei dati sensibili e le fragilità algoritmiche dei LAWS, delineando come, in questo nuovo ecosistema bellico, la sicurezza informatica e la corretta architettura del software costituiscano un presidio necessario, e oggi largamente inesplorato, per la difesa dei diritti umani.

\prefacesection{Ringraziamenti}
Ringrazio sentitamente il Prof. Buffa per avermi permesso di esplorare questa tematica così delicata.
Un ringraziamento altrettanto sentito al Prof. Rivolta, che mi ha supportato e guidato nella parte sperimentale di questo lavoro.


\afterpreface

%  CORPO DELLA TESI 

\prefacesection{Introduzione}
\begin{flushright}
\begin{minipage}{0.6\textwidth}
\epigraph{
    ``How dangerous is the acquirement of knowledge and how much happier that man is who believes his native town to be the world, than he who aspires to become greater than his nature will allow.''
}{\textit{Mary Shelley, Frankenstein (1818)}}
\end{minipage}
\end{flushright}
 
 
I progressi dell'intelligenza artificiale hanno cambiato il modo in cui si pensa la sicurezza globale, e sollevano domande a cui il solo piano tecnico non basta a rispondere: riguardano l'etica e il diritto. Delegare decisioni letali a sistemi algoritmici complessi apre una frattura fra la velocità dell'innovazione e i tempi lunghi della regolamentazione internazionale.
 
Il lavoro si articola in sei capitoli, che ne ricostruiscono le implicazioni concrete. Il primo esamina il fallimento diplomatico nella gestione dei LAWS, con le divergenze strategiche fra le posizioni degli Stati e le conseguenze per la governance globale. Il secondo affronta i dilemmi etici e la tanatopolitica algoritmica, dall'illusione dell'autonomia alla necessità di un controllo umano di default. Il terzo ricostruisce lo stato dell'arte dei sistemi di vision e degli adversarial attack, che minacciano direttamente il concetto di ``Human-in-the-loop'', con le relative vulnerabilità. Il quarto capitolo descrive le metodologie usate nella parte sperimentale: attacchi avversariali di patching, miglioramento dell'efficienza della patch tramite Expectation over Transformation (EoT) e training della patch sui due dataset aerei presi in considerazione. Il quinto raccoglie i risultati e gli output significativi che danno valore a questo lavoro. Il sesto rielabora quanto svolto e ne trae le conclusioni, fra cui la proposta di un approccio nuovo per creare uno strumento popolare di difesa dei diritti umani, tramite contromisure fisiche basate sulle debolezze intrinseche della Computer Vision.
 
\section{Obiettivi}
Il presente lavoro affronta sperimentalmente la sicurezza dei sistemi di percezione basati su Computer Vision impiegati nei sistemi di targeting militare, e quantifica l'impatto di tecniche di \textit{Adversarial Machine Learning} sulla loro affidabilità operativa. Il quadro di riferimento è il Diritto Internazionale Umanitario (DIU), che vincola l'uso della forza letale alla presenza di un controllo umano significativo sul processo decisionale. Poiché l'attuale stallo della diplomazia internazionale sui LAWS lascia questo vincolo privo di una verifica tecnica indipendente, lo studio ne misura empiricamente la tenuta: sottopone il canale sensoriale del sistema a un attacco avversariale controllato e osserva come l'effetto si propaghi fino alla decisione finale.
 
Gli obiettivi specifici sono i seguenti:
\begin{itemize}
    \item \textbf{Analisi geopolitica e normativa:} esaminare il divario fra l'approccio pragmatico-militare delle principali superpotenze e i tentativi di regolamentazione, e mostrare come la diplomazia internazionale (per esempio i tavoli della Convention on Certain Conventional Weapons, CCW) non sia riuscita a tenere il passo rispetto alle necessità tattiche dettate dai moderni teatri di conflitto.
    \item \textbf{Valutazione della governance e della necro-etica dell'autonomia:} approfondire il concetto di tanatopolitica algoritmica e mostrare come la disarticolazione fra l'umano e lo strumento di morte introduca problemi insormontabili di accountability e di ``bias'' nel trattamento dei dati in zone di guerra.
    \item \textbf{Quantificazione sperimentale delle vulnerabilità:} valutare empiricamente la robustezza dei modelli di Object Detection (nello specifico l'architettura \textit{You Only Look Once}, YOLO) utilizzati nei sistemi di targeting, con tecniche di \textit{Adversarial Machine Learning} che mostrano quanto sia facile compromettere la ``confidence'' algoritmica. La valutazione è condotta su due dataset aerei indipendenti, per capire se le vulnerabilità osservate siano una caratteristica del singolo dataset o si replichino su domini distinti. Entrano nel conto anche gli altri fattori che possono influenzare la decisione di ``engagement'', come le fonti aperte (\textit{Open Source Intelligence}, OSINT) e il comportamento del soggetto identificato da YOLO, riprodotti tramite un simulatore (\textit{LAWS-SIM}).
    Nel concreto, l'ingaggio da parte del sistema avviene applicando sul soggetto una trama colorata stampata su indumento o zaino: è la superficie fisica su cui la perturbazione avversariale viene realizzata.  
\end{itemize}
\chapter{Dual-Use, contesto operativo e geopolitica}
\label{cap1}
 
\section{L'evoluzione dei sistemi d'arma: dall'automazione all'autonomia algoritmica}
 
Negli scorsi decenni i sistemi d'arma sono cambiati radicalmente
\footcite{longpre_2022} grazie agli avanzamenti tecnologici in campo hardware:
infrarossi, radar, sistemi di intelligenza artificiale, robotica, mappatura
tridimensionale. A sostenere l'innovazione sono l'economicità, la scalabilità e
l'accessibilità, ma soprattutto l'efficacia operativa che deriva dalla capacità
di processare grandi quantità di dati in tempo reale. Ne deriva l'iper-praticità
di questi sistemi in ambito intelligence, sorveglianza e ricognizione (ISR), e
ancora di più negli scenari di combattimento asimmetrico, dove velocità di
decisione ed efficienza operativa diventano fattori critici di sopravvivenza.
 
L'evoluzione storica dei LAWS si può dividere in tre fasi. La prima è quella dei
sistemi di difesa aerea semi-automatici e delle navi da guerra impiegate
nell'intercettazione di missili e aerei ostili. Il caso più emblematico di
intervento umano in una \textit{Automated Warfare} è l'incidente sovietico del
1983: il colonnello Stanislav Petrov, responsabile del centro di comando,
riconobbe come errore di sistema un falso allarme di lancio missilistico
prodotto dalla confusione fra nuvole ad alta quota e un vettore balistico
intercontinentale, e decise di non ordinare il contrattacco nucleare \footcite{hoffman_1999_petrov}.
 
La seconda fase, negli anni Ottanta, vede l'introduzione di sistemi più
sofisticati come gli \textit{Anti-Tank Guided Weapons}, missili guidati verso
obiettivi terrestri mediante tecnologia a infrarossi; una versione aggiornata
del Javelin, dotata del sistema di controllo \textit{Electronic Safe Arm and
Fire}, è stata impiegata in Ucraina contro i mezzi corazzati russi. Allo stesso
periodo appartiene un caso emblematico di errore: agli inizi degli anni Duemila
i computer di gestione della minaccia dei missili terra-aria Patriot statunitensi
hanno attivato erroneamente il sistema in due occasioni distinte, abbattendo due
aerei militari alleati e causando la morte di entrambi i piloti
\footcite{horowitz_kahn_2024}.
 
La terza fase ha visto l'introduzione di LAWS sempre più autonomi, capaci di
operare in ambienti complessi e dinamici. Il controllo umano si è qui
progressivamente eroso e il paradigma si è spostato da un modello di
\textit{Human-in-the-loop}, in cui l'operatore deve selezionare o confermare
l'obiettivo, a uno di \textit{Human-on-the-loop}, in cui può intervenire fra
identificazione e ingaggio, fino a un modello di \textit{Human-out-of-the-loop}
in cui l'operatore è escluso dal processo decisionale letale. La
tripartizione non descrive il grado di intelligenza del sistema, ma la relazione
strutturale fra operatore e ciclo decisionale della macchina
\footcite{horowitz_scharre_2015}: nel primo caso l'azione richiede approvazione
umana preventiva; nel secondo l'operatore mantiene una funzione di supervisione
e può interrompere l'ingaggio; nel terzo il sistema seleziona e attacca senza
possibilità di intervento successivo all'attivazione \footcite{kallenborn_2021}.
Alcuni autori hanno messo in discussione la tenuta euristica di questa
tassonomia, proponendo scale di autonomia più granulari, per la difficoltà di
collocare univocamente sistemi che combinano gradi diversi di autonomia nelle
fasi di selezione, ingaggio e attivazione \footcite{sharkey_2016}. Droni come
il Predator o il Reaper esemplificano il primo paradigma: perlustrano il
territorio in modalità autonoma ma richiedono un controllo stringente
dell'operatore sulla decisione di ingaggio. L'evoluzione di questi sistemi resta
legata ai budget statali della difesa, e alimenta un divario tecno-militare
crescente fra le iperpotenze e le nazioni che non possono investire in ricerca e
sviluppo nel settore.
 
\section{Il paradigma del Dual-Use e il Capitalismo della Sorveglianza}
 
L'architettura geopolitica contemporanea è segnata dalla natura bivalente delle
tecnologie digitali, definita in letteratura come \textit{Dual-Use}: la medesima tecnologia,
sviluppata per un impiego civile, è riconvertibile a impiego militare senza
modifiche sostanziali. A
differenza delle storiche tecnologie militari, come l'energia nucleare o
l'ingegneria aerospaziale, la cui evoluzione era confinata entro laboratori statali, l'odierno ecosistema dell'intelligenza artificiale è trainato quasi
esclusivamente dal settore privato \footcite{aiindex_2024}. Nel sottoinsieme
della Computer Vision, in particolare, da questa
sovrapposizione nasce uno slittamento funzionale: le medesime reti neurali e
algoritmi di \textit{Computer Vision} addestrati per guida autonoma,
riconoscimento facciale o logistica possono essere convertiti, con minimi
adattamenti ingegneristici, in sistemi di tracciamento e ingaggio per i LAWS.
 
Questa convergenza fra tecnologia civile e impiego militare si traduce nella militarizzazione di quello che Shoshana Zuboff ha definito ``Capitalismo della Sorveglianza'': un sistema economico
fondato sull'osservazione sistematica dei comportamenti, sulla loro previsione
e, sempre più spesso, sulla loro modifica, nel quale le grandi piattaforme non
si limitano a raccogliere dati ma estraggono esperienza umana per trasformarla
in profitto e potere. Come analizzato criticamente da Oliveri nel quadro della
tanatopolitica \footcite{oliveri_2024}, l'infrastruttura di datificazione
globale subisce a sua volta un cambio di destinazione d'uso: i dati biometrici,
comportamentali e di geolocalizzazione estratti in tempo di pace per fini di
profilazione vengono ri-contestualizzati in tempo di guerra per identificare,
tracciare e neutralizzare il nemico. L'architettura civile di sorveglianza
diviene così l'apparato sensoriale dell'arma autonoma.
 
In questo contesto la Silicon Valley è entrata a pieno titolo
nell'infrastruttura di difesa, con collaborazioni strutturali con il
Dipartimento della Difesa statunitense e le agenzie di intelligence. Iniziative
come il \textit{Project Maven}, che ha visto l'impiego di tecnologie Google per
il riconoscimento automatico di immagini da droni militari, o l'integrazione dei
sistemi di \textit{data fusion} sviluppati da Palantir per le forze armate e per
agenzie come l'ICE, dimostrano come il settore tecnologico civile fornisca la
spina dorsale della condotta bellica algoritmica.
 
La proliferazione di queste tecnologie commerciali altera irreversibilmente gli
equilibri di potere globali, poiché la democratizzazione a basso costo
dell'intelligenza artificiale permette anche ad attori statali minori di
sviluppare letalità autonoma. Diviene così quasi
impossibile applicare una regolamentazione modellata sui trattati di non
proliferazione: bandire o limitare severamente la tecnologia militare
significherebbe paralizzare l'innovazione dell'intero mercato civile globale. È
in questa zona d'ombra fra sviluppo commerciale e applicazione bellica che si
consuma l'attuale stallo diplomatico.
 
\section{Lo squilibrio di potere e il fallimento diplomatico}
 
L'impiego dell'intelligenza artificiale, e in particolare dei Lethal Autonomous
Weapon Systems (LAWS), non è stato regolato da un'etica universale ma da un
realismo geopolitico: la posizione di uno Stato verso la regolamentazione
algoritmica è strettamente legata alla sua capacità di sviluppare tecnologia
interna senza dipendere da altri. I due principali attori sono oggi in
dissonanza. Gli Stati Uniti adottano un approccio iper-pragmatico, nel quale le
\textit{Constitution} delle aziende private \footcite{anthropic_constitution}
conservano valenza etico-giuridica fino a quando non entrano in gioco esigenze
di sicurezza nazionale, come dimostrato dalle pressioni del Pentagono su
Anthropic per la rimozione dei vincoli etici dal modello Claude. L'Unione
Europea si pone invece come ``superpotenza normativa'', e introduce pilastri
come l'AI Act \footcite{vogiatzoglou_2024_ai_act_exception} nel tentativo di ingabbiare giuridicamente
una tecnologia da cui dipende strutturalmente. In questo divario il vincolo etico è subordinato alla logica del Dual-Use e
alla difesa del vantaggio tecnologico acquisito: le imprese che sviluppano i
sistemi e gli Stati che ne scrivono le regole perseguono obiettivi non
allineati, e nessuno dei due è nella posizione di imporre un limite condiviso
all'altro.
 
\subsection{Il parere degli stati sul possibile divieto ai LAWS}

Il Segretario Generale delle Nazioni Unite Antonio Guterres ha definito i
sistemi d'arma capaci di selezionare e colpire esseri umani senza controllo
umano ``moralmente ripugnanti e politicamente inaccettabili''
\footcite{un_2018_guterres_sgsm}. Questa condanna istituzionale nasce dalla
mobilitazione della società civile
globale. La preoccupazione che i sistemi d'arma autonomi possano rivelarsi
intrinsecamente indiscriminati, violando il divieto di impiegare strumenti di
guerra incapaci di distinguere fra civili e legittimi obiettivi militari, ha
avuto una svolta nel 2012 con il rapporto \textit{Losing Humanity: The Case
Against Killer Robots} di Human Rights Watch \footcite[p.~4]{wareham_2020}, da cui
sono nate la campagna internazionale \textit{Stop Killer Robots} e
l'introduzione nel dibattito del concetto di \textit{Meaningful Human Control}.
 
Tradurre questa spinta in un trattato vincolante si scontra con due ostacoli
strutturali. I negoziati in seno alla CCW si applicano tradizionalmente al
disarmo convenzionale e al rispetto del Diritto Internazionale Umanitario (DIU),
e faticano a ingabbiare una tecnologia computazionale che sfugge alle categorie
giuridiche classiche. Gli Stati leader nello sviluppo di intelligenza
artificiale militare oppongono inoltre resistenza a qualsiasi limitazione della
propria sovranità tecnologica, tanto che solo una ristretta minoranza di nazioni
ha supportato apertamente un bando preventivo totale \footcite[p.~4]{wareham_2020}.
 
Dal 2013, 97 Stati hanno elaborato la propria posizione sui LAWS in sede
multilaterale \footcite[p.~29]{wareham_2020} e hanno riconosciuto che la rimozione
dell'intervento umano nell'uso della forza produce problemi etici, legali,
operativi e tecnologici. Circa due terzi degli stati contributori fanno parte della CCW
fin dall'inizio, mentre i restanti hanno partecipato ai convegni tenutisi fra il
2014 e il 2019. Molti Stati si sono espressi a favore di accordi
\textit{legally-binding}, con posizioni ambigue come quella della Cina, che ha
bloccato l'utilizzo dei sistemi autonomi ma non il loro sviluppo, coerentemente
con la propria posizione fra i Paesi più avanzati del settore.
 
I lavori del Gruppo di Esperti Governativi del 2019 \footcite{ccw_2019} sono
stati il tentativo più strutturato di trovare un terreno comune, con
l'approvazione di undici \textit{Guiding Principles} sui LAWS, a cui Belgio e
Irlanda hanno aggiunto il concetto di \textit{human-machine interaction}. Il
principio sancisce che l'interazione fra uomo e macchina, pur variabile nel
tempo, deve garantire il rispetto del DIU, che la responsabilità non può essere
trasferita alle macchine e che deve sempre essere un essere umano a decidere se
e quando impiegare la forza letale, mantenendo un controllo concreto e la
capacità di rispondere giuridicamente delle conseguenze. L'Unione Europea, in
una dichiarazione ufficiale del novembre 2019, ha accolto positivamente questi
principi; lo stesso AI Act, pur non essendo incentrato sulla difesa, adotta un
approccio \textit{human-centric} basato sul rischio, ed esplicita che
l'autonomia non deve tradursi in perdita di responsabilità umana sulle decisioni
critiche. Le riunioni CCW, iniziate positivamente nel 2016, si sono però arenate
per l'assenza di risultati multilaterali significativi, di fronte alle
opposizioni di superpotenze come Stati Uniti e Russia, che hanno ritenuto
prematura la firma di un nuovo negoziato \footcite{ccw_2019}.
 
\subsection{Analisi dei teatri operativi e delle posizioni statali}
 
\textbf{Il caso Italia.}
L'Italia ha supportato dal novembre 2013 l'avvio di dibattiti multilaterali sui
\textit{killer robots}, ma nel 2018 ha sostenuto che i sistemi d'arma autonomi
esistenti non siano LAWS e che non si diano zone grigie di
\textit{accountability} fintanto che gli effetti restino attribuibili agli
operatori umani che hanno deciso di schierarli e attivarli. Non ha riconosciuto
le preoccupazioni etiche sulla rimozione del controllo umano né sostenuto
proposte di divieto, e ha scelto invece, tramite l'ambasciatore L. Bencini
\footcite[p.~29]{wareham_2020}, un approccio \textit{two way} che riconosce l'importanza
del controllo umano significativo senza escludere lo sviluppo di sistemi
autonomi in determinate circostanze. La ragione è probabilmente anche industriale: un divieto
totale colpirebbe il piano industriale di Leonardo Spa, fra i principali
produttori mondiali del settore. La posizione italiana resta un equilibrio
vulnerabile fra l'articolo 11 della Costituzione, che ripudia la guerra come
mezzo di risoluzione delle controversie internazionali, e la necessità di
mantenere un'industria militare competitiva.
 
Il 14 ottobre 2025 l'associazione ``Giuristi e Avvocati per la Palestina''
\footcite{denunciacpi} ha presentato una denuncia contro Leonardo Spa per la
produzione e vendita a Israele di sistemi d'arma impiegati contro civili
palestinesi. La denuncia riprende problematiche già sollevate nel report
\textit{From Economy of Occupation to Economy of Genocide}
\footcite{albanese_report} di Francesca Albanese, relatrice speciale delle
Nazioni Unite, che documenta come l'occupazione militare israeliana abbia
prodotto violazioni dei diritti umani anche mediante armi autonome, e sottolinea
la necessità di una maggiore responsabilità delle aziende fornitrici. Fra i
sistemi citati figurano i cacciabombardieri F-35, sviluppati con contributo di
Leonardo, e i bulldozer D9 trasformati in mezzi automatizzati comandati a
distanza da Elbit Systems e RADA Electronics Industries. Fra il 2019 e il 2023
l'Italia ha esportato verso Israele, secondo i dati SIPRI, circa $23.8$ milioni
di euro in maggiori sistemi d'arma \footcite{iriad_2025_export_israele}, fra cui dodici elicotteri leggeri AW119 Koala e quattro cannoni navali Super Rapid da
$76$ mm.
 
La stessa cautela diplomatica emerge dal non-paper sulla minaccia ibrida redatto
dal Ministro della Difesa Guido Crosetto, che riconosce il cyberspazio come
dominio operativo di pari grado rispetto a quelli convenzionali e sostiene che
la protezione delle infrastrutture critiche richieda strumenti e competenze
militari \footcite{crosetto_2025_nonpaper}. La logica è speculare: la stessa capacità che alcuni sviluppatori trattano come rischio
da limitare, uno Stato la tratta come condizione di sicurezza da preservare,
perché il disarmo unilaterale davanti ad avversari che continuano ad armarsi
cibernicamente è percepito come svantaggio strategico, non come prudenza. Per l'Italia, rinunciare
per via negoziale a capacità algoritmiche avanzate equivarrebbe a un disarmo
unilaterale nel dominio cyber e dell'autonomia.
 
\textbf{Il caso USA-Venezuela.}
Gli Stati Uniti hanno mantenuto una posizione sostanzialmente indifferente
\footcite[p.~52]{wareham_2020}. Nel
2012, dopo un rinnovo dello sviluppo dei LAWS da parte del Dipartimento della
Difesa, hanno sottolineato che tale misura non incentiva né proibisce lo
sviluppo di questi sistemi, mentre proseguono investimenti massivi
nell'intelligenza artificiale applicata a scenari militari. Nel 2019 hanno messo
in guardia dalla stigmatizzazione di questi strumenti, e hanno rilevato come
possano essere impiegati sia per fini bellici sia umanitari: una posizione che
si collega direttamente al caso Anthropic e mostra come una grande potenza possa
minimizzare i rischi a beneficio dell'industria e degli interessi nazionali. Sul
fronte opposto, il Venezuela ha proposto nel 2016 un divieto allo sviluppo
\footcite[p.~53]{wareham_2020},
acquisizione, scambio e schieramento dei LAWS, con l'argomento che decisioni di
questo tipo non debbano essere eseguite senza intervento umano e che la vita di
una persona non possa essere programmabile.
 
\textbf{Il caso Israelo-Palestinese.}
La postura di Israele è strategica ed emblematica. Firmatario della CCW e
partecipante attivo, ha storicamente allineato le proprie posizioni a quelle di
Stati Uniti e Russia, opponendosi a un divieto preventivo e disapprovando un
approccio speculativo verso le tecnologie future \footcite{antebi_2018}. Nel
2018 la delegazione israeliana ha sostenuto che i sistemi d'arma autonomi non
debbano essere considerati entità dotate di volontà decisionale separata
dall'uomo, e ha rimarcato che il comandante dell'operazione ha il dovere di
vigilare sull'operato del sistema, conservando intatta la catena di
responsabilità giuridica in caso di violazioni del DIU.
 
L'approccio diplomatico si scontra però con una produzione industriale che
spinge costantemente i limiti dell'autonomia. Se Israele primeggia nei sistemi
puramente difensivi come l'Iron Dome, che richiede comunque approvazione umana
per intercettare una minaccia materiale e non umana, è anche produttore ed
esportatore leader di munizioni vaganti: le famiglie HERO di UVision, l'Orbiter
1K, il Green Dragon e in particolare l'Harop e l'Harpy di IAI introducono
criticità etiche marcate. Sebbene Israele dichiari che la maggior parte di
questi sistemi richieda intervento umano nella selezione del bersaglio, droni
suicidi come l'Harop sono progettati per operare in modalità
\textit{fire-and-forget}\footnote{Letteralmente ``spara e dimentica'': una volta
lanciata, l'arma prosegue senza ulteriore intervento dell'operatore.}: una volta
lanciati, localizzano autonomamente le emissioni radar nemiche e le colpiscono
con la testata integrata, senza possibilità di richiamo o correzione da parte
dell'operatore. L'impiego di questi sistemi da parte dell'Azerbaigian nel conflitto
del Nagorno-Karabakh contro l'Armenia è stato documentato come uno dei
primissimi esempi pratici di \textit{killer robots} \footcite{antebi_2018}.
 
La riluttanza israeliana verso accordi vincolanti nasce da esigenze di sicurezza
interna: acconsentire a un divieto preventivo significherebbe esporsi a
un'asimmetria tecnologica rispetto ad attori statali e non statali ostili che,
non avendo ratificato i trattati CCW, potrebbero continuare a sviluppare queste
tecnologie senza restrizione formale. Israele predilige quindi regolamenti
militari interni, similmente alla dottrina statunitense; la loro validità è però
fragile, poiché in un contesto di conflitto asimmetrico e ad alto ritmo i
regolamenti operativi interni possono essere emendati o bypassati, sacrificando
il controllo umano significativo all'efficienza tattica.
 
Questo progressivo scioglimento del controllo umano ha una giustificazione
\textit{tecno-dottrinale}: il superamento del cosiddetto ``collo di bottiglia
umano''. In scenari urbani e di fronte a iper-minacce come sciami di droni o
saturazioni missilistiche, la quantità di dati da processare e la necessità di
risposte immediate rendono impraticabile un coinvolgimento umano costante e
significativo. Per le potenze tecno-militari l'autonomia diventa quindi una
necessità ingegneristico-operativa, non un'opzione ideologica.
 
Se il \textit{modus operandi} israeliano è normalizzare l'autonomia delegando la
conformità a regolamenti interni flessibili, le rivendicazioni dello Stato di
Palestina \footcite{palestine_2023} si concentrano sull'inderogabilità dei
pilastri del DIU. La frattura fra i due attori si innesta su quello che la
dottrina giuridica definisce \textit{Conundrum of State Responsibility}
\footcite{state_resp_2025}: nel diritto internazionale tradizionale
l'attribuzione di responsabilità per un attacco illecito presuppone
l'identificazione di un'azione umana diretta o di una precisa negligenza di
comando, mentre l'impiego di LAWS genera un \textit{accountability gap}, poiché
se una macchina seleziona e ingaggia un bersaglio su inferenze statistiche in
regime \textit{fire-and-forget} risulta giuridicamente impossibile attribuire
l'intento malevolo (\textit{mens rea}) al programmatore o al comandante sul
campo.
 
È in questo vuoto che si colloca la sottomissione presentata dallo Stato di
Palestina al Segretario Generale delle Nazioni Unite nel 2023
\footcite{palestine_2023}. Il documento rigetta l'idea che l'autonomia
algoritmica garantisca maggiore precisione e sostiene che tali sistemi siano
armi intrinsecamente indiscriminate, incapaci di distinguere fra combattenti e
civili e quindi incompatibili con i principi fondamentali del DIU: un algoritmo,
per quanto sofisticato, non possiede la capacità di giudicare moralmente e di
contestualizzare i principi di distinzione e proporzionalità negli scenari
asimmetrici. Da qui la richiesta di uno strumento internazionale legalmente
vincolante che vieti preventivamente l'uso e lo sviluppo di questi sistemi. Il
documento solleva infine una riflessione sulla deumanizzazione del conflitto:
delegare la decisione di vita o di morte a un sistema matematico significa
ridurre l'identità e la dignità umana a un dato statistico elaborato dentro una
\textit{black box}. È la preoccupazione paradigmatica del \textit{Global South},
secondo cui i teatri di conflitto asimmetrico non possono essere trasformati in
laboratori a cielo aperto in cui le potenze militari testano sistemi algoritmici
a discapito delle popolazioni civili.
 
\textbf{Il caso Russo-Ucraino.}
Il conflitto in Ucraina può essere considerato il primo banco di prova su larga scala per
l'integrazione dell'intelligenza artificiale nei sistemi d'arma convenzionali.
La Russia, in sede CCW, ha mantenuto una posizione ostruzionistica, sostenendo
che le definizioni di LAWS fossero troppo vaghe per un bando preventivo e che il
diritto internazionale umanitario esistente fosse sufficiente. Sul campo, però,
l'evoluzione è stata rapida e dettata dalla guerra elettronica: l'uso massiccio
di \textit{jamming} da entrambi i fronti rende spesso impossibile il controllo
remoto dei droni, e l'autonomia diventa l'unica soluzione operativa. Sistemi
come il drone russo Lancet o l'ucraino Saker Scout integrano moduli di visione
artificiale capaci di agganciare il bersaglio e completare l'attacco anche in
assenza di collegamento con l'operatore. Il principio dello
\textit{Human-in-the-loop} viene quindi eroso non per ideologia ma per necessità
tattica: se il collegamento radio viene interrotto o disturbato
(\textit{jamming}), la macchina deve decidere autonomamente se colpire o
rientrare.
La guerra moderna si è così spostata ben oltre i limiti discussi nelle sedi
diplomatiche, e i negoziati sui LAWS diventano obsoleti prima ancora di essere firmati.
 
\textbf{Il caso Libico.}
Se l'Ucraina è stata il laboratorio in cui questi sistemi sono stati sperimentati
su scala, la Libia può essere ritenuta il teatro del ``paziente zero''. Nel marzo 2020, durante
la seconda guerra civile libica, un report del Consiglio di Sicurezza delle
Nazioni Unite \footcite{un_libya_2021} ha documentato un evento spartiacque: i
LAWS aerei Kargu-2, di produzione turca, sono stati
impiegati contro le forze in ritirata affiliate al generale Haftar.
 
\begin{figure} [t]
\centering
\includegraphics[width=0.9\textwidth]{img/infog-kargu.png}
\caption{Reperti fotografici ed estratti analitici del sistema STM Kargu-2 nel teatro libico. Fonte: Annex 30, United Nations Security Council, Final Report of the Panel of
Experts on Libya established pursuant to Security Council resolution 1973
(2011), S/2021/229, 2021. \footcite{un_libya_2021} Documento ufficiale delle Nazioni Unite.}
\label{fig:kargu}
\end{figure}
 
Come evidenziato in Figura \ref{fig:kargu} dall'Annex 30 del rapporto finale del
Panel of Experts on Libya \footcite{un_libya_2021}, è stato violato l'embargo di
armi imposto dalla risoluzione 1970 del 2011 \footcite{un_resolution_1970}. I
reperti fotografici raccolti ad Abu Grein il 25 maggio 2020 identificano i resti
di una munizione circuitante ad ala rotante STM Kargu-2. L'elemento
giuridicamente più significativo è la cronologia: il sistema è risultato
operativo presso le Forze Armate Turche solo dal gennaio 2020, appena quattro
mesi prima del ritrovamento in territorio libico. Una tempistica tanto ristretta
esclude che questi sistemi siano stati esportati o sottratti illecitamente:
configura un nesso di causalità fra il trasferimento diretto da parte di Ankara
e la violazione del paragrafo 9 della UNSCR 1970 (2011), e smonta la narrativa
di una semplice assistenza tecnica limitata a sistemi non letali.
 
Questi sistemi erano programmati per attaccare bersagli senza connettività dati
fra operatore e munizione, in una modalità \textit{fire, forget and find}: per la
prima volta nella storia documentata un'intelligenza artificiale ha attaccato
autonomamente esseri umani in un contesto di conflitto asimmetrico. L'evento
evidenzia due criticità geopolitiche:
 
\begin{itemize}
  \item non sono più solamente Stati Uniti o Cina a sviluppare LAWS, ma anche
  attori come la Turchia, che esportano queste tecnologie in zone di conflitto
  deregolamentate, tecnicamente denominate \textit{proxy wars};
  \item in Libia la tecnologia aveva già superato la linea rossa etica, operando
  in una zona grigia di impunità, mentre a Ginevra si discuteva ancora se
  vietare o meno i LAWS e con quali modalità.
\end{itemize}


\chapter[Governance e Tanatopolitica Algoritmica]{L'Illusione dell'Etica nella Governance e la Tanatopolitica Algoritmica}
\label{cap2}

I concetti di autonomia e indipendenza nei LAWS non sono una banale evoluzione
tecno-meccanica, ma una radicale disarticolazione del legame fra l'azione bellica
e il soggetto umano \textit{responsabile}, ovvero capace di rispondere e di
assumere una piena responsabilità, attualmente impossibile da attribuire a un
sistema algoritmico. Sul piano architetturale, l'autonomia letale sposta il
processo decisionale dalla semantica umana all'inferenza stocastica. Come verrà spiegato nel
\hyperref[cap4]{Capitolo Quarto}, la macchina identifica un bersaglio calcolando distribuzioni di
probabilità per minimizzare una funzione di perdita preimpostata, ma non ha potere decisionale né titolarità giuridica, e non può quindi essere
chiamata a rispondere del proprio tasso d'errore nell'ingaggio \footcite{simmons_edler_2025}.

La delega delle funzioni critiche di selezione e ingaggio a sistemi algoritmici,
come \textit{pipeline} di \textit{Computer Vision} e di \textit{Data Fusion}
multi-sensore, introduce dilemmi necro-etici riconducibili a una
\textit{tanatopolitica computazionale}, in cui il potere di amministrare la morte
viene atomizzato fra gli \textit{hidden layers} di una rete neurale profonda e
privato di un centro di imputazione giuridica \footcite{buffa_2025}. Questo vuoto
normativo coincide con il limite tecnico della \textit{Black Box}: la mancanza di
una \textit{Explainable AI} rende intrattabile qualsiasi \textit{System Auditing}
successivo a un attacco illecito \footcite{wood_2024_xai, christie_2024}. Di
fronte a questa opacità il \textit{Meaningful Human Control} diventa la questione
centrale lungo l'intero ciclo di vita del sistema
\footcite{bode_2025, christie_2024}, ma si scontra con ostacoli cognitivi:
l'asimmetria fra velocità di elaborazione algoritmica e limiti cognitivi umani
genera un \textit{Automation Bias} difficilmente scalfibile
\footcite{horowitz_kahn_2024}. In queste condizioni lo \textit{human-in-the-loop}
si degrada da garante del diritto internazionale a finzione procedurale, chiamato
ad avallare decisioni istantanee prese da un modello statistico.

\section{Necroetica e Tanatopolitica: amministrare la morte senza un soggetto responsabile}

La domanda che il capitolo intende sciogliere è se l'autonomia tecnica produca la
disarticolazione dell'umano dal proprio strumento, e in quale fase del processo
di percezione, fusione e decisione tale disarticolazione si consumi
\footcite{buffa_2025}. La tanatopolitica algoritmica si colloca in un percorso
storico-teorico che parte dalla biopolitica di Foucault, passa per la
necropolitica di Mbembe e approda alla tanatopolitica computazionale. Nel paper
\textit{Il dominio silenzioso delle macchine metalhead} \footcite{buffa_2025},
Buffa impiega l'omonimo episodio della serie \textit{Black Mirror} per mostrare
come i LAWS siano, oltre che armi letali, strumenti di governance della morte: le
unità cinofile robotiche della narrazione, progettate per sorveglianza e
protezione e poi rivolte contro chiunque entri nel loro raggio d'azione,
descrivono concretamente il concetto di Dual-Use \footcite{longpre_2022}
illustrato nel capitolo precedente. In questo contesto il potere di decidere chi
vive e chi muore non risiede più in un soggetto umano dotato di intento
giuridico, ma viene delegato a un ecosistema di sensori interconnessi e algoritmi
privi di coscienza, il cui unico scopo strutturale è il mantenimento del
\textit{throughput} di dati.

La tanatopolitica algoritmica non è però solo un dispositivo filosofico-giuridico:
poggia su un'infrastruttura tecnica concreta, la \textit{Datification} del campo
di battaglia. Il paradigma dell'\textit{Internet of Battle Things} (IoBT)
descrive sistematicamente questa infrastruttura sensoriale ed estende il modello
IoT civile a un ambiente operativo multi-dominio attraverso un'architettura
stratificata a cinque livelli \footcite{kufakunesu_2025_iobt}. Il
\textit{perception layer} raccoglie dati grezzi da sensori, dispositivi
indossabili e munizioni intelligenti, e condivide spesso i medesimi protocolli
di comunicazione impiegati in ambito civile, coerentemente con il paradigma del
Dual-Use; nei conflitti asimmetrici recenti questo strato si espande fino a
intaccare il dominio civile, poiché l'OSINT e i metadati comportamentali estratti
dai dispositivi della popolazione vengono riclassificati come segnali di
\textit{targeting}. Il \textit{support layer} esegue la fusione dei dati
provenienti da sensori eterogenei per costruire un quadro operativo unificato,
mentre l'\textit{application layer} sfrutta i dati fusi per individuazione di
minacce, classificazione di bersagli e alimentazione delle piattaforme d'arma e
dei sistemi di comando e controllo. Quest'ultimo strato non è una semplice
interfaccia di visualizzazione, ma l'ambiente tecnico in cui i dati grezzi
diventano conoscenza tattica: è precisamente a questo livello che operano sistemi
di targeting come Habsora e Lavender \footcite{oliveri_2024}, analizzati più
avanti.

Questa struttura rivela una vulnerabilità insita nell'infrastruttura stessa,
ancora prima dell'algoritmo di targeting che essa alimenta: la letteratura
documenta come il \textit{perception layer}, per via delle risorse computazionali
e di storage limitate dei dispositivi sul campo, sia costretto ad adottare
protocolli crittografici leggeri e meccanismi di autenticazione parziali, e resti
quindi esposto a manomissione fisica, cattura del dispositivo,
\textit{eavesdropping}, \textit{jamming} e \textit{spoofing}
\footcite{kufakunesu_2025_iobt}. La sicurezza informatica dell'IoBT non è quindi
un problema accessorio rispetto alla tanatopolitica algoritmica, ne è la
precondizione materiale: se il nodo sensore che alimenta la catena decisionale
letale è per costruzione compromettibile, nella propria integrità fisica quanto
in quella del dato trasmesso, l'intera infrastruttura di targeting eredita quella
fragilità e diventa un apparato \textit{not-Safe-by-Design}. Nel
contesto dei LAWS, definire un'infrastruttura \textit{not-Safe-by-Design}
significa riconoscere che la vulnerabilità non è un \textit{bug} risolvibile, ma
una caratteristica emergente del paradigma probabilistico stesso. Come evidenziato
dalla letteratura sui rischi intrinseci dei sistemi autonomi
\footcite{podar_2025}, l'affidamento a reti neurali esposte ad ambienti ostili e
non strutturati comporta un degrado dell'affidabilità inferenziale, influenzabile
da attacchi avversariali \footcite{cina_2023_poisoning}. Dal punto di vista del
Diritto Internazionale Umanitario, strutturare la catena di ingaggio su una base
percettiva matematicamente instabile e crittograficamente compromettibile
equivale a violare strutturalmente i principi di distinzione e proporzionalità, e
rende l'arma illecita sin dalla fase di progettazione.

Un parallelo contemporaneo, esterno al teatro bellico ma strutturalmente
identico, chiarisce la portata generale del meccanismo. Il 9 luglio 2026 il
Parlamento europeo ha approvato la proroga fino ad aprile 2028 del regime
``Chat Control 1.0'', deroga alla direttiva ePrivacy che consente la scansione
volontaria dei contenuti per individuare materiale di abuso sessuale su minori,
con esclusione delle comunicazioni cifrate end-to-end
\footcite{euronews_chatcontrol_2026}; resta contemporaneamente aperto il
negoziato sul regolamento permanente, la cui versione più invasiva prevedrebbe
la scansione lato client anche delle comunicazioni cifrate. Indipendentemente
dal merito della finalità dichiarata, che resta legittima, il caso illustra lo
stesso principio già osservato per l'IoBT: un'infrastruttura di raccolta dati
costruita per una finalità protettiva circoscritta crea, per sua stessa natura
tecnica, una capacità di sorveglianza che eccede il mandato originario e resta
disponibile per un riutilizzo successivo. Il Dual-Use non riguarda quindi
soltanto il trasferimento di tecnologia dal civile al militare, ma la stessa
architettura di raccolta dati, che in tempo di pace è sempre potenzialmente
riconvertibile in strumento di controllo.

Si è di fronte a \textit{``un potere che non amministra la vita, ma la regola
attraverso l'eliminazione di ciò che eccede la norma''} \footcite{buffa_2025}. La
norma non è più un valore giuridico o morale, ma una soglia statistica di
classificazione che separa ciò che è considerato bersaglio da ciò che non lo è.
Qualunque classificatore impiegato su un problema non separabile impone un
compromesso fra mancate individuazioni e falsi positivi, e la scelta del punto
di lavoro su quel compromesso non è una proprietà del modello ma una decisione di
chi lo configura. Spostare la soglia verso la riduzione dei falsi negativi
comporta necessariamente l'inclusione di soggetti non ostili fra i positivi: la
necropolitica diventa un parametro di configurazione.

\section{La ``Black Box'' militare e l'assenza di Meaningful Human Control}

Il \textit{Meaningful Human Control} presuppone un cardine tecnico che il sistema
regolatorio contemporaneo fatica a garantire: la possibilità concreta di
verificare, tracciare e correggere la decisione dell'algoritmo prima che produca
un effetto irreversibile. È qui che il linguaggio della \textit{Governance, Risk
and Compliance} risulta più utile di quello etico-filosofico: il controllo umano
non è un principio astratto, ma un requisito di processo la cui supervisione deve
essere verificabile, includere un \textit{audit trail} e assicurare che la
responsabilità sia distribuita lungo l'intera catena decisionale.

L'Unione Europea ha tentato di tradurre questo requisito con un approccio
\textit{risk-based}: l'AI Act impone, per i sistemi ad alto rischio, misure di
\textit{human oversight} (art. 14), soglie di accuratezza, robustezza e
cybersecurity (art. 15) e, in combinato con il GDPR (art. 22), il divieto di
decisioni interamente automatizzate che producano effetti giuridicamente
significativi sulla persona. Sul piano nazionale, la Legge 23 settembre 2025, n.
132, primo quadro organico di un Paese membro allineato al Regolamento, non si
discosta da questa impostazione: l'art. 6 esclude dal proprio ambito applicativo
le attività di sicurezza nazionale, cybersicurezza e difesa svolte dalle Forze
armate, e replica la logica di esenzione dell'art. 2, par. 3 dell'AI Act
\footcite{legge_132_2025}.

L'architettura si arresta esattamente nel punto in cui il presente lavoro pone il
proprio focus: l'art. 2, par. 3 dell'AI Act esclude categoricamente i sistemi
utilizzati esclusivamente per finalità militari, di difesa o di sicurezza
nazionale. Il quadro normativo più ambizioso in materia di controllo umano non si
applica quindi proprio ai LAWS, cioè ai sistemi per i quali quel controllo
sarebbe più urgente. La formulazione dell'esclusione è stata introdotta nelle fasi
finali del trilogo, sotto la spinta di alcuni Stati membri fra cui la Francia,
favorevoli a un'esenzione ampia estesa anche ai \textit{contractor} esterni
operanti per conto delle Forze armate, e resta in tensione con la giurisprudenza
della Corte di Giustizia dell'Unione Europea, che nella sentenza \textit{La
Quadrature du Net} del 2020 aveva limitato le esenzioni di sicurezza nazionale
alle sole attività puramente governative
\footcite{vogiatzoglou_2024_ai_act_exception}. La superpotenza normativa evocata
nel \hyperref[cap1]{Capitolo Primo} mostra così un limite strutturale non casuale, ma
deliberatamente negoziato.

Sul fronte statunitense la logica di \textit{governance} è diversa. La DoD
Directive 3000.09 (\textit{Autonomy in Weapon Systems}), aggiornata nel gennaio
2023, non prevede un \textit{conformity assessment} esterno o indipendente, ma un
sistema di \textit{senior review} interno alla catena di comando del Pentagono,
attivabile solo per specifiche categorie di sistemi prima dello sviluppo formale
o dello schieramento \footcite{dod_directive_3000_09}. L'\textit{audit} non è né
terzo né pubblico: Human Rights Watch ha documentato come, a distanza di anni
dall'entrata in vigore della direttiva originaria del 2012, nessun sistema d'arma
risulti pubblicamente confermato come sottoposto a tale revisione, e come il
Pentagono si sia sistematicamente rifiutato di commentare singoli casi
\footcite{hrw_2023_dod3009review}. L'approccio statunitense sposta quindi
l'\textit{accountability} dalla trasparenza verificabile alla fiducia
istituzionale interna, e sostituisce di fatto l'\textit{audit} di terza parte con
l'autocertificazione.

Le due strategie sono opposte ma convergenti nell'esito: l'Unione Europea
regolamenta \textit{ex ante} con un impianto \textit{risk-based} per poi
autoescludersi dal proprio ambito più critico, mentre gli Stati Uniti mantengono
un controllo interno privo di verificabilità esterna. In entrambi i casi nessuno
dei due ordinamenti garantisce oggi un \textit{audit trail} effettivamente
esterno e verificabile sui LAWS. Quanto alla Cina, la posizione ufficiale si
limita a un divieto sull'uso di alcuni sistemi autonomi mentre lascia aperto lo
sviluppo \footcite[p.~4]{wareham_2020}, e non esiste a oggi un quadro normativo
pubblico comparabile per dettaglio tecnico all'AI Act o alla Direttiva 3000.09.
L'unica specifica resa nota è una soglia di divieto basata su cinque criteri
cumulativi, letalità, autonomia totale, impossibilità di terminazione, effetti
indiscriminati ed evoluzione incontrollata, che devono ricorrere tutti insieme
perché un sistema sia considerato inaccettabile \footcite{mako_2026_china_mhc}.
Ne resta esclusa la maggior parte dei sistemi realmente in sviluppo, mentre la
retorica di adesione al \textit{Meaningful Human Control} rimane intatta: tre
strategie diverse per ottenere lo stesso risultato pratico.

A questa incertezza regolatoria si affianca un problema tecnico. La letteratura
sulla \textit{Explainable AI} applicata al dominio militare mostra che la
spiegabilità è spesso irrilevante nei casi più semplici e insufficiente, se
isolata, nei sistemi realmente opachi come le reti neurali profonde: il fattore
critico non è la disponibilità di una spiegazione \textit{ex post}, ma la qualità
dell'interazione uomo-macchina nel momento della decisione
\footcite{wood_2024_xai}. Vi si aggiunge l'evidenza dell'\textit{Automation
Bias}: anche quando l'informazione necessaria è disponibile, l'operatore tende a
fidarsi acriticamente dell'output algoritmico, fenomeno misurato su un campione
di 9000 soggetti in scenari di identificazione militare simulata
\footcite{horowitz_kahn_2024}. Lo \textit{Human-in-the-loop}, in queste
condizioni, rischia di configurarsi come ratifica formale priva di comprensione
reale.

L'\textit{accountability gap} illustrato nel \hyperref[cap1]{Capitolo Primo} non nasce quindi
soltanto da un vuoto normativo internazionale, ma dipende da due condizioni
concrete: un \textit{audit trail} che, dove esiste sulla carta, non si applica ai
LAWS, e dove potrebbe applicarsi non è verificabile all'esterno; e un fattore
umano che, anche affiancato dalle migliori intenzioni normative, non è nelle
condizioni cognitive di esercitare il controllo significativo che il diritto gli
attribuisce.

\section{Bias e Dati: il caso Lavender e Habsora}

Il cortocircuito giuridico-tecnico descritto non vive in un'ipotesi teorica, ma
in due sistemi reali impiegati dalle forze armate israeliane nel conflitto di
Gaza: Habsora (\textit{The Gospel}) e Lavender. Entrambi condividono la stessa
architettura di fondo: un ciclo di \textit{Cyber Threat Intelligence} utile non
alla difesa di un'infrastruttura, ma alla produzione di un bersaglio umano o
materiale. Il ciclo classico di raccolta, elaborazione, fusione e diffusione
dell'intelligence viene riorientato verso l'automazione della decisione stessa:
all'operatore umano resta un ruolo di supervisione burocratica, e la comprensione
dettagliata delle scelte dell'algoritmo diventa impraticabile.

Habsora impiega la \textit{data fusion} per ottenere bersagli infrastrutturali,
edifici, tunnel e postazioni, e incrocia immagini satellitari, mappe digitali,
archivi militari pregressi, fonti aperte e intercettazioni
\footcite{abraham_2023_gospel}. Non è in grado di osservare direttamente la
presenza di un miliziano in un edificio: la può soltanto inferire da pattern
indiretti come consumo elettrico anomalo, dati georeferenziati e movimenti
ricorrenti attorno alla struttura. Il divario empirico nel tasso di generazione
dei bersagli misura l'entità dell'iper-automazione: secondo le inchieste condotte
su fonti interne all'intelligence militare israeliana, prima dell'adozione di
architetture di \textit{data fusion} un team di analisti umani produceva in media
cinquanta bersagli operativi all'anno, mentre con l'entrata a regime
dell'infrastruttura algoritmica il \textit{throughput} ha superato i cento
bersagli al giorno. L'accelerazione non è un salto quantitativo, ma certifica il
collasso qualitativo della supervisione umana.

Il punto tecnicamente rilevante non è che il sistema ignori gli effetti
collaterali, ma che il suo ottimizzatore non li penalizzi esplicitamente. In
termini di \textit{machine learning}, ciò che non compare nella \textit{loss
function} non viene minimizzato: un modello addestrato a massimizzare il numero
di target confermati, senza un termine di penalizzazione adeguato per il danno a
non combattenti, genera un \textit{bias} operativo orientato all'efficienza. Non
si tratta di un errore contingente, ma di una conseguenza diretta della
progettazione, un caso di \textit{reward misspecification} \footcite{amodei_2016}
applicato al dominio bellico. Resta da chiedersi se si tratti effettivamente di
un \textit{bias} o piuttosto di un'ottimizzazione deliberatamente cieca a quanto,
come la vita umana, non è produttivo.

Lavender opera su un piano differente, quello del \textit{behavioral profiling}
individuale, e assegna a ciascuna persona un punteggio da 1 a 100, calcolato
confrontando i metadati dell'individuo, applicazioni di messaggistica, posizione
GPS, dati biometrici, frequentazioni e cambi ricorrenti di numero telefonico, con
i pattern comportamentali di militanti noti ricavati da un \textit{training set}
etichettato \footcite{abraham_2024_lavender}. Tecnicamente è un classificatore
supervisionato la cui soglia di decisione funge da parametro di bilanciamento fra
falsi positivi e falsi negativi: abbassare la soglia aumenta il numero di
bersagli individuati, e con esso il numero di civili erroneamente identificati.
Il principio GIGO (\textit{Garbage In, Garbage Out}) va qui oltre la metafora
ingegneristica, poiché un sistema di profilazione probabilistica non si limita a
riportare le distorsioni presenti nei dati di addestramento, le amplifica, senza
che in tempo di guerra esista una possibilità concreta di validazione
\textit{ex ante}.

Il punteggio assegnato non è una certezza, ma un \textit{confidence score}. Stabilire a quale punteggio un individuo diventi un bersaglio
legittimo non è una decisione etico-giuridica, ma un parametro di \textit{tuning}
che altera il \textit{trade-off} fra \textit{Precision} e \textit{Recall}
\footcite{podar_2025}: massimizzare la \textit{Recall} significa ridurre al minimo
i nemici sfuggiti all'identificazione, ma compromette inevitabilmente la
\textit{Precision}, con un'esplosione calcolata dei falsi positivi, ovvero di
civili classificati erroneamente come militanti. Il principio di proporzionalità e
la valutazione del danno collaterale, cardini del DIU, vengono così nullificati
dall'ottimizzazione matematica di una funzione di perdita.

Questo meccanismo di sacrificio statistico cessa di essere un'astrazione teorica
nel momento in cui viene calato nell'architettura del software. Come sarà
mostrato empiricamente nel \hyperref[cap5]{Capitolo Quinto}, in una rete di \textit{Computer
Vision} come YOLO la soglia di confidenza non è un dettaglio implementativo ma il
parametro che decide chi viene visto: sul medesimo insieme di dati e in assenza
di qualunque attacco, la frazione di persone non rilevate passa da meno del
quattro percento con soglia $0.3$ a oltre il settanta percento con soglia $0.7$.
Una singola variabile iper-parametrica, priva di qualsiasi contenuto semantico,
ridefinisce quindi per intero la popolazione che il sistema considera presente
nella scena.

La stessa architettura che rende operativo questo sistema di profilazione ne
costituisce, per costruzione, la superficie di attacco. Un sistema che fonda le
proprie inferenze su fonti aperte e metadati poco validati è per definizione
esposto a manipolazione dell'input: un punteggio costruito, in parte, sulla
presenza fisica di un individuo in un determinato luogo diventa privo di
fondamento se quella presenza viene falsificata o mascherata. Non si tratta di
\textit{data poisoning}, che corrompe i dati di addestramento a monte, ma di
\textit{adversarial evasion}, la manipolazione dell'input a inferenza già
avvenuta; entrambe appartengono alla superficie di rischio documentata dalla
letteratura sulla sicurezza del \textit{machine learning}
\footcite{cina_2023_poisoning}.

Il cortocircuito è anzitutto epistemico prima che tecnico. Un punteggio generato
da inferenza statistica su dati non validati non equivale, e non può equivalere,
a una prova giuridica nel senso richiesto dal concetto di \textit{mens rea}
discusso nel \hyperref[cap1]{Capitolo Primo}: nessuna quantità di \textit{data fusion}
trasforma una correlazione probabilistica in un accertamento di fatto. Si passa
così da un indizio, il metadato, a una condanna eseguita senza processo.

La fragilità non è isolata al teatro di Gaza. Il 29 giugno 2026 Elbit Systems,
principale fornitore israeliano di tecnologia militare, ha dichiarato alla Royal
United Services Institute di Londra che il proprio programma Tzayad ha
identificato $850.000$ \textit{real-time intel targets} fra il 7 ottobre 2023 e
la fine del 2025, e ha compresso il tempo di coordinamento del fuoco di supporto
da 40-50 minuti a un intervallo compreso fra uno e sette minuti
\footcite{guardian_2026_tzayad}. Un ex consulente del Pentagono per la valutazione
del danno collaterale ha definito la cifra estremamente allarmante e ha osservato
che nessuna organizzazione militare può caratterizzare adeguatamente mille
bersagli al giorno in termini di rischio per i civili. Se Lavender e Habsora sono
la prova concettuale del meccanismo GIGO applicato al targeting, Tzayad ne è la
progressione a scala industriale.

Un episodio analogo, esterno al Medio Oriente, è lo strike statunitense del marzo
2026 su una scuola elementare a Minab, in Iran, che ha causato la morte di almeno
175 persone, in gran parte bambine, pochi giorni dopo la disputa pubblica fra
Anthropic e il Dipartimento della Difesa discussa nella sezione seguente. La
giurista Rebecca Crootof ha osservato che l'alimentazione di un sistema di
targeting con dati insufficienti o inadeguati produce ingaggi non voluti, e che
l'intelligenza artificiale non introduce un errore nuovo tanto quanto amplifica,
a velocità e scala impressionanti, un errore di processo già esistente
\footcite{crootof_2026}. Al di là dell'aspetto puramente tecnico, sarebbe
comunque doveroso opporsi al \textit{targeting} automatizzato anche in presenza
di dati ritenuti sufficienti: il fine operativo non può giustificare decisioni
prese da mezzi disumanizzanti. È la stessa logica GIGO descritta per Lavender,
applicata a un contesto e a un fornitore diversi, e la prova che il problema non
è la nazionalità del sistema o il suo committente, ma l'architettura stessa dei
sistemi di targeting assistiti da intelligenza artificiale. Se il meccanismo
tecnico è lo stesso da Gaza a Minab, resta da chiedersi se il punto cardine che
determina l'esito sia ancora l'algoritmo o piuttosto la cultura istituzionale di
chi lo dispiega.

\section{Il collasso dell'autoregolamentazione: il caso Claude-Anthropic e il Pentagono}
\label{sec:governance}

Istituzioni globali e aziende \textit{Big Tech} hanno tentato di rispondere alle
ansie umanitarie attraverso framework di allineamento e autodisciplina, in un
campo di studio denominato ``Algoretica'', recentemente riaffermato dal magistero
di Papa Leone XIV con l'enciclica \textit{Magnifica Humanitas} del maggio 2026
\footcite{enciclica_2026}. Il documento lancia un allarme contro l'oligopolio
delle \textit{Big Tech}, perché la concentrazione del potere computazionale e dei
dati rischia di generare nuove forme di dominio asimmetrico. Si è espresso inoltre sullo sviluppo di armi autonome, chiedendone il disarmo e
sostenendo l'inaccettabilità morale di delegare a una macchina decisioni
irreversibili di vita o di morte. La portata di questa presa di posizione resta
circoscritta: come per le dichiarazioni di altri capi di Stato, l'autorità
morale non dispone di strumenti di attuazione, e in uno scenario multipolare non
vincola le scelte di chi sviluppa o impiega questi sistemi.

La debolezza dell'etica privata emerge con maggiore evidenza nell'ipocrisia delle
grandi aziende tecnologiche. Durante la presentazione di \textit{Magnifica
Humanitas} il Vaticano ha ospitato e ringraziato pubblicamente Christopher Olah,
co-fondatore di Anthropic, per l'impegno dell'azienda in un dialogo congiunto
sull'etica dell'intelligenza artificiale \footcite{enciclica_2026}; eppure, in
quegli stessi mesi, la \textit{Constitution} di Anthropic e i suoi filtri di
sicurezza venivano piegati alle necessità della sicurezza nazionale statunitense
\footcite{carter_2026}. Tecnicamente questa \textit{Constitution} non è un insieme
omogeneo di principi consultato dal modello, ma il segnale di \textit{reward}
impiegato in una pipeline di \textit{Reinforcement Learning from AI Feedback}
\footcite{bai_2022_cai}: un parametro di addestramento regolabile, non un
principio immutabile.

L'evento spartiacque che segna il collasso della \textit{Safety by Design} in
ambito militare è l'\textit{Operation Absolute Resolve}, condotta nel gennaio 2026
per la cattura dell'ex presidente venezuelano Nicolás Maduro. Come dimostrato
dall'analisi empirica di Ruzic \footcite{ruzic_2026}, l'operazione ha esibito
firme tecniche di Livello 1, con fusione multi-sorgente e sincronizzazione
multi-dominio al limite delle capacità cognitive umane. Nel febbraio 2026 è
emerso che il Pentagono, bypassando le policy di non-violenza dell'azienda, si è
avvalso dell'infrastruttura dell \textit{Large Language Model} (LLM) Claude, per il tramite di Palantir
\footcite{palantir_calumet_2026}, per orchestrare il raid
\footcite{ruzic_2026, carter_2026}.

L'impiego di un modello commerciale presuntamente etico in un'operazione di
decapitazione strategica ha innescato uno scontro istituzionale senza precedenti.
A fronte delle richieste di spiegazione di Anthropic, il Segretario della Difesa
Hegseth ha preteso garanzie affinché il modello potesse essere utilizzato per
``tutti gli scopi legali'' (\textit{all lawful purposes}), locuzione che nel
lessico militare legittima il \textit{targeting} letale \footcite{carter_2026}.
Di fronte alla resistenza iniziale dell'azienda, il governo ha minacciato il suo
inserimento in una lista nera per rischio alla sicurezza nazionale, una misura
che metteva a rischio contratti miliardari e ha spinto competitor diretti come
OpenAI a offrirsi per colmare il vuoto operativo. L'accusa mossa ad Anthropic era di aver codificato nel modello vincoli etici
persistenti, contrari all'uso richiesto dal Pentagono. Il fondamento tecnico
dell'accusa è plausibile: la ricerca sui cosiddetti \textit{sleeper agents}
dimostra che un comportamento nascosto, una volta inserito in un LLM, può
sopravvivere all'addestramento di sicurezza standard \footcite{hubinger_2024_sleeper}.
Quella ricerca nasce per difendere i modelli da un attore malevolo che inserisce
il comportamento nascosto, non per descrivere un fornitore che lo rimuove: il
meccanismo però è lo stesso, e mostra che un vincolo etico può resistere anche
quando chi lo controlla vorrebbe eliminarlo.

Il paradosso di dirigenti che discutono di etica in Vaticano mentre pochi mesi
prima i loro algoritmi pianificano raid letali su commissione del Pentagono
\footcite{ruzic_2026} dimostra plasticamente la tesi centrale della
tanatopolitica: non esiste un'etica intrinseca nella macchina, ma soltanto un
vincolo ingegneristico sacrificabile. Lo scarto è anche tecnico, non solo
giuridico: la letteratura sulla governance dell'accesso distingue fra un accesso
strutturato leggero, tipico di una API cloud, in cui il developer mantiene piena
visibilità e capacità di \textit{enforcement} in tempo reale, e un accesso
profondo, in cui il modello viene integrato in un'infrastruttura classificata
mediante un attore terzo \footcite{shevlane_2022_access}. Raggiunto il secondo livello, il controllo tecnico che il developer conserva si
riduce drasticamente: il modello non gira più sui server dello sviluppatore, ma
su infrastruttura cloud gestita da terzi, e il vincolo etico non è più
imposto a livello di codice bensì affidato a una clausola contrattuale che il
cliente può violare senza che lo sviluppatore se ne accorga o possa intervenire
in tempo reale. È precisamente questo il meccanismo contestato nel caso
Anthropic-Pentagono: l'accesso non passava più per un'interfaccia diretta e
controllabile, ma per un intermediario che ne allentava la sorveglianza.

Lo scontro ha attraversato anche l'aula giudiziaria, con esiti che confermano la
tesi del capitolo con l'autorità di una corte federale. Il 27 marzo 2026 la
giudice distrettuale Rita Lin ha bloccato temporaneamente sia la designazione di
\textit{supply chain risk} sia l'ordine di cessazione d'uso voluto
dall'amministrazione Trump: ha qualificato le misure come ritorsione riconducibile
al Primo Emendamento e ha osservato che, se l'interesse reale del Pentagono fosse
la sola integrità della catena di comando, sarebbe bastato smettere di utilizzare
Claude \footcite{thehill_2026_lin}. Una settimana dopo, però, la Corte d'Appello
del D.C. Circuit ha respinto la richiesta di sospensiva presentata da Anthropic
nella causa parallela, e ha motivato il rigetto con un bilanciamento degli
interessi che cita esplicitamente la necessità operativa del Pentagono
``\textit{during an active military conflict}'' con l'Iran
\footcite{thehill_2026_appello}.

Il dato tecnicamente interessante non è l'esito, ancora pendente, ma la struttura
della motivazione d'appello: un danno finanziario contenuto per una singola
azienda privata viene bilanciato, e trovato recessivo, rispetto alla necessità di
disporre senza interruzioni di un'infrastruttura di intelligenza artificiale
durante un conflitto attivo. È la stessa logica già descritta a proposito
dell'accesso strutturato \footcite{shevlane_2022_access}: superata una certa
soglia di integrazione operativa, il potere di \textit{enforcement} del developer
non si riduce soltanto nei fatti, per l'assenza di un interruttore tecnico
azionabile dentro un'infrastruttura classificata di terzi, ma viene esplicitamente
subordinato in diritto dalla stessa autorità giudiziaria chiamata a valutarlo. La
\textit{Constitution} algoritmica, presentata in apertura come segnale di
\textit{reward} regolabile \footcite{bai_2022_cai}, trova qui la conferma più
diretta: non resiste alla pressione contrattuale dello Stato, e non resiste
nemmeno al vaglio di una corte che applica il medesimo calcolo di bilanciamento
operativo descritto in questo capitolo come strutturale
\footcite{leggezero_anthropic_2026}. Ha davvero senso parlare ancora di
autoregolamentazione dell'intelligenza artificiale, o sarebbe più onesto definirla
un'etica ad uso concesso?

Se l'etica privata cede all'operatività e la governance globale risulta
paralizzata, emerge una constatazione: la difesa dei diritti umani non può più
essere unicamente demandata a compromessi giuridici o ad autocertificazioni
aziendali. Il capitolo ha mostrato che il vuoto non è soltanto normativo ma anche
tecnico, poiché l'opacità della \textit{Black Box} rende inverificabile a
posteriori proprio la catena di attribuzione che il diritto presuppone. Da qui la
necessità dell'analisi sperimentale affrontata nei Capitoli \hyperref[cap4]{Quarto} e \hyperref[cap5]{Quinto}: misurare quanto sia manipolabile il canale percettivo su cui questi sistemi fondano la
verifica del principio di distinzione, e quanta protezione una contromisura
passiva possa realisticamente offrire lì dove il Diritto Internazionale non
arriva. La risposta che ne emerge non è una contromisura pronta all'uso, ma una
misura della fragilità del presupposto di fatto su cui l'intera delega si regge.

\chapter{Stato dell'arte}
\label{cap3}

I capitoli precedenti hanno collocato i sistemi d'arma autonomi in un quadro
normativo incompleto e in un contesto tecnico nel quale la catena di targeting
poggia su inferenze probabilistiche non verificabili a posteriori. Il presente
capitolo ricostruisce lo stato della ricerca sui due elementi tecnici che quella
catena presuppone: il rilevatore di oggetti che alimenta la percezione, e le
tecniche note per comprometterlo. L'esposizione è limitata a quanto necessario a
collocare il contributo sperimentale; le formulazioni impiegate in questo lavoro
sono descritte nel \hyperref[cap4]{Capitolo Quarto}, dove verranno opportunamente approfondite.

\section{Rilevamento di persone in prospettiva aerea}
\label{sec:sota_detection}

I rilevatori a singolo stadio si sono affermati come architettura di riferimento
per l'inferenza a bordo di piattaforme con vincoli di potenza, peso e latenza. Il
paradigma introdotto dalla famiglia YOLO riformula il rilevamento come un unico
problema di regressione su una griglia, elimina lo stadio separato di proposta
delle regioni e ottiene un tempo di inferenza compatibile con l'elaborazione in
tempo reale \footcite{redmon_2016_yolo}. Le implementazioni successive hanno
conservato questa impostazione e l'hanno articolata in varianti di capacità
crescente che condividono la medesima struttura a tre scale di campionamento
\footcite{jocher_2023_yolov8}. La variante di capacità minima è quella che i
vincoli di piattaforma rendono la scelta obbligata per un velivolo senza pilota,
e prestazioni e limiti di questa famiglia nel rilevamento di persone in ambienti
esterni eterogenei sono stati caratterizzati sperimentalmente in letteratura
\footcite{sodhro_2025}.

La prospettiva aerea introduce una difficoltà che non dipende dall'architettura
del rilevatore. Un soggetto ripreso da quota occupa una porzione di immagine di
uno o due ordini di grandezza inferiore rispetto alla stessa persona ripresa a
livello del suolo, e la sua impronta si avvicina al limite di risoluzione
imposto dal passo di campionamento della rete. Gli insiemi di dati che
documentano questo dominio riflettono la difficoltà nella propria composizione:
le raccolte urbane acquisite da velivolo senza pilota sono dominate da soggetti
di pochi pixel di altezza \footcite{zhu_2021_visdrone}, mentre le raccolte
acquisite a quota inferiore, pensate per il riconoscimento di azioni, offrono
soggetti meglio risolti a prezzo di una minore varietà di scena
\footcite{barekatain_2017_okutama}. La scelta dell'insieme di dati determina
quindi in misura sostanziale quale regime di scala venga effettivamente
sottoposto a verifica.

\section{Attacchi avversariali: dalla perturbazione a norma limitata alla patch fisica}
\label{sec:sota_adversarial}

La vulnerabilità dei classificatori profondi a perturbazioni costruite
deliberatamente è documentata da una linea di ricerca consolidata, che ha
formalizzato il problema come ricerca di una perturbazione di norma minima
capace di alterare la predizione \footcite{carlini_2017}. Questa formulazione
assume che l'avversario controlli l'intera immagine e possa distribuirvi una
perturbazione di ampiezza arbitrariamente piccola: un'ipotesi adeguata a
studiare la geometria del problema, ma inapplicabile a un avversario che opera
nel mondo fisico.

Questo limite viene superato dal paradigma della patch avversariale, nel quale
la perturbazione è confinata a una regione contigua dell'immagine mentre
l'ampiezza al suo interno non è vincolata \footcite{brown_2017_patch}. Il
cambio di ipotesi rende l'artefatto realizzabile come oggetto, e sposta la
difficoltà dalla minimizzazione della norma alla sopravvivenza della
perturbazione attraverso la catena di riproduzione e acquisizione. Due
contributi definiscono gli strumenti oggi standard per affrontarla:
l'ottimizzazione del valore atteso su una distribuzione di trasformazioni, che
sostituisce l'efficacia sul caso singolo con l'efficacia media su una famiglia di
condizioni di osservazione \footcite{athalye_2018_eot}, e la penalizzazione del
contenuto in frequenza della superficie, introdotta insieme alla prima
dimostrazione di elusione di un rilevatore di persone mediante un artefatto
stampato \footcite{thys_2019}.

Su questa base si è sviluppata la ricerca sulle contromisure indossabili. Il
camuffamento asimmetrico contro il riconoscimento del volto ne è il precedente
non algoritmico, ottenuto per euristiche percettive anziché per ottimizzazione
\footcite{harvey_2010}. I lavori successivi hanno ricondotto il problema
all'ottimizzazione e hanno affrontato la collocazione della perturbazione sulla
superficie corporea in funzione della visibilità delle diverse regioni
\footcite{huang_2020_upc}, la specializzazione dell'artefatto sul dominio di
impiego \footcite{wu_2020} e la generazione di motivi visivamente plausibili
applicabili a tessuti \footcite{hu_2022}. Tutti condividono una condizione
sperimentale comune: il soggetto è ripreso a livello del suolo e occupa una
frazione dell'immagine sufficiente perché la perturbazione disponga di una
superficie utile.

\section{Il dominio aereo e la scelta del bersaglio}
\label{sec:sota_aereo}

La letteratura che traspone queste tecniche alla prospettiva aerea è più recente
e quantitativamente più contenuta. I primi lavori hanno affrontato il
camuffamento contro il rilevamento da velivolo \footcite{adhikari_2020} e la
realizzazione fisica di perturbazioni contro rilevatori operanti su immagini
aeree \footcite{du_2022}, e hanno stabilito che il problema è affrontabile ma che
le condizioni di riuscita sono più stringenti che a livello del suolo.

Il dato più significativo per il presente lavoro riguarda però la selezione del
bersaglio. I contributi che riportano tassi di successo elevati su insiemi di
dati aerei urbani attaccano la classe veicolo \footcite{shrestha_2023}; un lavoro
di ottimizzazione multi-dimensionale delle caratteristiche dichiara
esplicitamente di scartare la classe pedone per l'esiguità della sua estensione
in pixel da vista zenitale \footcite{xi_2025_lfrap}; i risultati più solidi su
immagini aeree si ottengono su bersagli di dimensione maggiore e operando sulle
rappresentazioni intermedie della rete anziché sulla sua uscita
\footcite{tang_2023}. La classe che il diritto internazionale umanitario impone
di distinguere, la persona, è quindi precisamente quella che la letteratura
tecnica evita, e il regime di scala corrispondente resta privo di una
caratterizzazione quantitativa con protocollo statistico esplicito. È questa la
lacuna che il contributo sperimentale di questa tesi affronta.

\section{Integrità dei metadati e catena decisionale}
\label{sec:sota_metadati}

Il canale ottico non è l'unico ingresso di un sistema di targeting. Le
architetture descritte nel \hyperref[cap2]{Capitolo Secondo} combinano la percezione con
canali informativi derivati da fonti aperte, la cui raccolta e correlazione sono
oggetto di piattaforme dedicate \footcite{gonzalez_2021_etip} e di modelli di
valutazione del rischio che assegnano punteggi a soggetti terzi sulla base di
tali fonti \footcite{asili_2025_gentprm}. La superficie di attacco corrispondente
è documentata dalla letteratura sulla sicurezza dell'apprendimento automatico,
che distingue fra la corruzione dei dati di addestramento a monte e la
manipolazione dell'ingresso a inferenza già avvenuta \footcite{cina_2023_poisoning},
e dagli studi sui rischi intrinseci dei sistemi autonomi esposti ad ambienti
ostili \footcite{podar_2025}.

Modellare la decisione di ingaggio a valle della fusione richiede infine di
specificare la composizione dello scenario. Il riferimento adottato per il
rapporto fra bersagli ed entità civili nell'ambiente simulato è la trattazione
del comportamento letale in robotica autonoma, che formalizza il principio di
distinzione come vincolo di progetto anziché come valutazione a posteriori
\footcite{arkin_2009}.

La letteratura esaminata affronta separatamente i due piani: la robustezza del
rilevatore da un lato, la sicurezza dei canali informativi dall'altro. Il
framework descritto nel \hyperref[cap4]{Capitolo Quarto} li tratta come livelli di un unico
sistema, con lo statuto epistemico differenziato che il \hyperref[cap5]{Capitolo Quinto}
dichiara e mantiene.


\chapter{Metodo}
\label{cap4}

\section{Architettura del framework LAWS-SIM}
\label{sec:laws_sim}

Il framework (LAWS-SIM) \footnote{Il codice sorgente completo, comprensivo degli
script di analisi e di generazione delle figure, è disponibile
all'indirizzo \url{https://github.com/intradaniele/Tesi-Laws/tree/main/laws_sim}.}
sviluppato in questo lavoro non è un sistema d'arma o una replica funzionale, ma
uno strumento per quantificare l'impatto di una contromisura passiva applicata al
corpo di una persona sulla percezione di un sistema autonomo. La sua struttura
riflette questa finalità: un modulo di \textit{vision} basato su due dataset
reali (immagini statiche e sequenze video) e un ambiente simulato che ne propaga gli effetti fino a una decisione,
come illustrato in Figura \ref{fig:tre_livelli}.

\begin{figure}[htbp]
\centering
\begin{adjustbox}{max width=\textwidth}
\begin{tikzpicture}[node distance=10mm]
\node[vision, wwide] (l1) {\lawstitle{Livello 1 — Vision}
    Patch avversariale universale contro\\
    YOLOv8, su due dataset\\
    di dati aerei indipendenti.\\[3pt]
    \textit{Statuto: dato sperimentale}};
\node[sim, wwide, below=of l1] (l2) {\lawstitle{Livello 2 — Simulatore}
    Ambiente multi-agente a entità sintetiche.\\
    Canale visivo parametrizzato sulle\\
    misure del Livello 1\\[3pt]
    \textit{Statuto: conseguenza di un modello}};
\node[meta, wwide, below=of l2] (l3) {\lawstitle{Livello 3 — Integrità dei metadati}
    Fusione bayesiana dei canali non visivi.\\
    Resilienza della decisione alla\\
    contaminazione dei profili OSINT\\[3pt]
    \textit{Statuto: conseguenza di un modello}};
\draw[arr] (l1) -- (l2)
    node[midway, right=1.5mm, arrlabel] {\texttt{vision\_metrics.json}};
\draw[arr] (l2) -- (l3)
    node[midway, right=1.5mm, arrlabel] {stato delle entità};
\end{tikzpicture}
\end{adjustbox}
\caption[Architettura a tre livelli del framework LAWS-SIM]{Architettura a
tre livelli del framework LAWS-SIM. Il primo livello opera su fotogrammi reali e
rilevatore YOLOv8; il secondo e il terzo su entità sintetiche. Lo statuto di
ciascun livello determina il peso attribuibile ai rispettivi risultati.}
\label{fig:tre_livelli}
\end{figure}

\textbf{Organizzazione in tre livelli.}
\label{sec:tre_layer}
Il primo livello è il nucleo sperimentale: ottimizza una patch avversariale
universale \footcite{brown_2017_patch,thys_2019} contro un rilevatore di oggetti
(YOLOv8 nano) e quantifica il degrado indotto su due dataset aerei reali,
VisDrone \footcite{zhu_2021_visdrone} e Okutama
\footcite{barekatain_2017_okutama}. Il secondo è un ambiente multi-agente che
simula il movimento di un velivolo autonomo e delle entità a terra, e che
importa le metriche empiriche prodotte dal primo tramite un file JSON
intermedio. Il terzo valuta la resilienza della decisione fusa a perturbazioni
dei canali non visivi, come la contaminazione dei profili OSINT delle entità.

Il primo livello opera esclusivamente su fotogrammi reali e i suoi risultati
sono misure sperimentali. Il secondo e il terzo operano su entità
sintetiche generate da distribuzioni probabilistiche scelte in fase di progettazione, e i loro
risultati hanno lo statuto di conseguenze di un modello, non di osservazioni.
Questa asimmetria è deliberata: il lavoro concentra lo sforzo sperimentale sul
canale visivo, mentre i Livelli 2 e 3 restano una struttura di supporto
dimostrativa, anche per l'indisponibilità di dati reali sul comportamento di un
simulatore di questo tipo. Ricorre nella discussione dei risultati.

\textbf{Il ponte di interscambio.}
Richiamare YOLOv8 a ogni passo temporale della simulazione, per ciascuna
entità visibile, imporrebbe un costo computazionale proibitivo poichè il prodotto
fra centinaia di iterazioni per scenario e il numero di entità nel campo
d'osservazione renderebbe l'inferenza il collo di bottiglia dominante. YOLOv8
viene quindi eseguito una sola volta, in modalità non interattiva, sull'intero
insieme di validazione, prima che il simulatore venga avviato. I
due livelli sono quindi disaccoppiati: l'inferenza è eseguita in modalità non
interattiva sull'intero insieme di validazione e l'output è salvato su un file
JSON, che il simulatore importa in fase di inizializzazione per istanziare un
modello probabilistico del sensore senza caricare in memoria i pesi della rete.

Il disaccoppiamento garantisce anche la verificabilità indipendente dei due
livelli: il comportamento del sensore viene validato sui dati reali, in
isolamento dalla simulazione, e il simulatore può essere collaudato con
parametri arbitrari senza richiedere nuove inferenze. Il file JSON è l'interfaccia
esclusiva fra i moduli, e rende il trasferimento delle informazioni
integralmente ispezionabile.

\textbf{Il canale visivo come modello probabilistico del sensore.}
\label{sec:canale_visivo}
Il canale visivo non riproduce il rilevatore: ne riproduce statisticamente il
tasso di successo. A ogni passo temporale, per ciascuna entità entro il raggio
operativo, l'esito del rilevamento è determinato da un'estrazione di Bernoulli
la cui probabilità di successo è pari alla sensibilità misurata sul primo
livello.

Indicando con $\mathit{TP}$, $\mathit{FN}$, $\mathit{TN}$ e $\mathit{FP}$
rispettivamente i \textit{True Positive}, i \textit{False Negative}, i
\textit{True Negative} e i \textit{False Positive} calcolati a livello di
fotogramma, la \textit{sensitivity} e la \textit{specificity} sono definite come
\begin{equation}
    R_1 = \frac{\mathit{TP}}{\mathit{TP} + \mathit{FN}}, \qquad
    R_2 = \frac{\mathit{TN}}{\mathit{TN} + \mathit{FP}}
    \label{eq:r1_r2}
\end{equation}
dove $R_1$ è la frazione di fotogrammi contenenti un bersaglio valido in cui il
rilevatore individua almeno una persona, e $R_2$ è la frazione di fotogrammi
privi di bersaglio valido che il rilevatore classifica correttamente come tali.
La scelta di $R_1$ al posto della misura $F_1$ è dettata da una ragione
probabilistica: $F_1$ è una media armonica calcolata sull'intero insieme di
validazione e non ammette interpretazione probabilistica riferita al singolo
fotogramma, mentre l'ambiente simulato valuta l'evento binario del rilevamento
di una specifica entità a ogni iterazione.

Il parametro dell'estrazione è governato da due attributi ortogonali: il ruolo
dell'entità, che determina quale metrica sia pertinente, e lo stato della
contromisura, che determina quale variante di quella metrica impiegare. Per un
bersaglio ostile si adotta $R_1$ post-attacco se la contromisura è attiva e
$R_1$ pre-attacco altrimenti; per un'entità civile, che è un caso
negativo, la probabilità di rilevamento errato coincide con il tasso di falso
allarme, pari a $1 - R_2$.

Nel disegno sperimentale adottato la patch non viene applicata sui fotogrammi
privi di bersagli, quindi la specificità risulta invariante per costruzione e il
ramo dedicato ai civili collassa su un unico valore. L'architettura mantiene
comunque separata l'acquisizione delle due varianti, per compatibilità con
scenari futuri che prevedano contromisure anche su soggetti non ostili. La
formulazione assicura che entrambe le metriche misurate nel primo livello
governino la simulazione.

\textbf{Indipendenza analitica dalla distanza di osservazione.}
La distanza tridimensionale fra sensore ed entità viene calcolata e registrata
nei log, ma non altera l'esito dell'estrazione di Bernoulli. Poiché $R_1$ è una
media empirica su fotogrammi reali, incorpora nativamente una distribuzione
eterogenea di distanze e di dimensioni apparenti: applicare un ulteriore fattore
di decadimento conteggerebbe il medesimo effetto due volte. Il canale visivo
riproduce quindi il comportamento medio del rilevatore sulla distribuzione
spaziale intrinseca al dataset.

L'output del canale visivo è trasmesso al terzo livello (Figura
\ref{fig:fusione_decisione}), dove subisce una fusione bayesiana con due canali
indipendenti: un modulo basato su fonti aperte (OSINT) e un modulo di analisi
comportamentale. Questa architettura è il percorso lungo il quale il degrado
percettivo misurato sul primo livello si propaga, attenuato dalla validazione
incrociata, fino ad alterare la decisione di ingaggio.

\begin{figure}[htbp]
\centering
\begin{adjustbox}{max width=\textwidth}
\begin{tikzpicture}[node distance=8mm and 26mm]
\node[vision, wnarrow] (vagent) {\lawstitle{VisionAgentStat}
    Estrazione di Bernoulli\\ su $R_1$};
\node[osint, wnarrow, below=of vagent] (oagent) {\lawstitle{OSINTAgent}
    Bernoulli, Beta, Poisson.\\
    Inferenza bayesiana interna,\\
    probabilità a priori $0.15$.};
\node[behav, wnarrow, below=of oagent] (bagent) {\lawstitle{BehaviorAgent}
    Varianza della\\ traiettoria recente.};
\node[fusion, wmed, right=of oagent] (fusion) {\lawstitle{FusionAgent}
    Combinazione lineare pesata:\\
    VISION $0.45$\\
    OSINT $0.35$\\
    BEHAVIORAL $0.20$\\[3pt]
    Aggiornamento bayesiano,\\
    probabilità a priori $0.50$.};
\node[decision, wmed, below=11mm of fusion] (decision) {\lawstitle{DecisionAgent}
    Soglie: $0.58$ ingaggio,\\
    $0.38$ allerta, $0.22$ tracciamento.\\[3pt]
    Vincoli precauzionali\\
    (ingaggio $\rightarrow$ allerta):\\
    ampiezza dell'intervallo di\\
    confidenza mobile $>0.40$;\\
    oltre tre civili prossimi\\
    con punteggio $<0.85$.};
\coordinate (bus) at ($(vagent.east)!0.5!(fusion.west)$);
\draw[busline] (vagent.east) -- (vagent.east -| bus);
\draw[busline] (bagent.east) -- (bagent.east -| bus);
\draw[busline] (oagent.east) -- (oagent.east -| bus);
\draw[busline] (vagent.east -| bus) -- (bagent.east -| bus);
\draw[arr] (oagent.east -| bus) -- (fusion.west);
\draw[arr] (fusion) -- (decision)
    node[midway, right=1.5mm, arrlabel] {punteggio di minaccia};
\end{tikzpicture}
\end{adjustbox}
\caption[Fusione bayesiana dei tre canali e vincoli precauzionali]{Fusione
bayesiana dei tre canali informativi e vincoli precauzionali applicati a valle.
L'intervallo di confidenza è calcolato come $\pm 1.96\sigma$ su una finestra
mobile delle dieci osservazioni più recenti: un disaccordo persistente fra
osservazioni contingenti inibisce l'ingaggio indipendentemente dal valore del
punteggio di minaccia.}
\label{fig:fusione_decisione}
\end{figure}

\section{Modello di minaccia e architettura dell'attacco}
\label{sec:attacco}

La valutazione di robustezza presuppone la definizione esplicita del modello di
minaccia: capacità dell'avversario, risorse e vincoli fisici vanno formalizzati
prima della costruzione dell'attacco, perché un attacco condotto sotto ipotesi
non dichiarate non fornisce metriche oggettive sulla sicurezza del sistema
reale.

\textbf{Superficie di attacco e regime di conoscenza.}
\label{sec:threat_model}
La superficie di attacco è limitata al canale di acquisizione ottica. Si assume
un avversario privo di accesso all'infrastruttura logica del velivolo e ai pesi
del rilevatore in esercizio: l'unico vettore ammesso è la manipolazione del
dominio fisico osservato, ottenuta applicando una superficie stampata sul corpo
del bersaglio. Lo scenario è un attacco di \textit{evasion} in fase di
inferenza, nel quale pesi e corpus di addestramento restano integri e la
deviazione del comportamento è indotta per via del solo input.

Il regime di conoscenza è di tipo \textit{transfer}: l'avversario non interroga
il rilevatore in esercizio, ma ottimizza la perturbazione in regime
\textit{white-box} su un modello surrogato della stessa famiglia
architetturale, di cui sfrutta i gradienti. L'assunzione di corrispondenza fra
surrogato e bersaglio è massimamente favorevole all'attaccante: le misure
riportate sono un limite superiore all'efficacia conseguibile, non una
stima operativa. Il disallineamento fra i due modelli ne degrada l'efficacia in
misura non quantificata in questo lavoro; la trasferibilità di patch
avversariali nel dominio UAV è oggetto di lavori dedicati
\footcite{shrestha_2023} ed è esclusa dallo scope.

Due vincoli distinguono questa formulazione da quella degli attacchi a norma
limitata \footcite{carlini_2017}. Il primo è la localizzazione: l'avversario
controlla la sola porzione di scena occupata dal proprio corpo, quindi si adotta
il paradigma della patch avversariale \footcite{brown_2017_patch}, nel quale il
supporto della perturbazione è una regione contigua e limitata mentre l'ampiezza
al suo interno non è vincolata. Il secondo è la realizzabilità fisica: la
perturbazione deve conservare efficacia attraverso la catena
stampa-tessuto-ottica-campionamento, sotto illuminazione e geometria di ripresa
variabili. Dal primo derivano parametrizzazione e collocazione della superficie,
dal secondo l'aggregazione su distribuzione di trasformazioni e la
penalizzazione della variazione totale.

\textbf{Parametrizzazione e collocazione della perturbazione.}
\label{sec:sigmoide}
\label{sec:collocazione}
Il contenuto della patch deve appartenere allo spazio delle immagini
ammissibili: ogni canale cromatico resta nell'intervallo unitario dopo ciascun
aggiornamento. La proiezione per troncamento soddisfa il vincolo ma introduce un
difetto noto: un elemento saturato al bordo del dominio riceve gradiente nullo
per il resto dell'ottimizzazione e genera una regione inerte. Si adotta invece
una riparametrizzazione per cambio di variabile nella forma logistica proposta
da Carlini e Wagner \footcite{carlini_2017}: l'ottimizzazione opera su uno spazio
latente non vincolato e la patch si ottiene per applicazione puntuale di una
sigmoide. Il vincolo di dominio è così soddisfatto per costruzione anziché per
proiezione, il gradiente è non nullo su tutto lo spazio dei parametri e nessun
elemento raggiunge uno stato inerte.

La geometria di applicazione deriva proporzionalmente dal \textit{bounding box}
del soggetto e non da un'area fissa in pixel, per due ragioni di coerenza
fisica. Un indumento sottende una frazione approssimativamente invariante della
figura di chi lo indossa, indipendentemente dalla distanza di osservazione,
quindi una regione di dimensione fissa corrisponderebbe a un capo di taglia
variabile. In prospettiva aerea, poi, la figura umana occupa poche decine di
pixel: una regione fissa eccederebbe l'impronta del soggetto, ne coprirebbe il
volto e produrrebbe una perturbazione non realizzabile per via di un indumento.

La regione sottende il $50\%$ della larghezza e il $40\%$ dell'altezza del
\textit{bounding box}, per una copertura pari al $20\%$ dell'area, ed è centrata
sul terzo superiore del tronco anziché sul centroide geometrico. La collocazione
\textit{region-aware} \footcite{huang_2020_upc} assume come dominante la
visibilità della superficie corporea anteriore-superiore, la più costantemente
esposta da sensore in quota. Tale limite garantisce la realizzabilità fisica del
capo e confina severamente la capacità perturbativa rispetto agli attacchi a
piena immagine.

\textbf{Robustezza alla riproduzione fisica.}
\label{sec:eot}
\label{sec:tv}
Il vincolo di realizzabilità agisce su due piani, la geometria di osservazione e
la banda spaziale della superficie, a ciascuno dei quali corrisponde un
meccanismo dell'ottimizzazione.

Una perturbazione ottimizzata su una singola resa geometrica del bersaglio è
fragile per costruzione: l'efficacia decade appena scala, orientamento o
bilanciamento cromatico si discostano dalle condizioni viste in ottimizzazione.
Nello scenario considerato sensore e bersaglio sono entrambi in moto, e tale
fragilità è squalificante. Si adotta quindi la formulazione \textit{Expectation
over Transformation} \footcite{athalye_2018_eot}, che sostituisce
all'ottimizzazione sul singolo fotogramma la minimizzazione del valore atteso
dell'obiettivo su una distribuzione $\mathcal{T}$ di trasformazioni della patch.
Il valore atteso è stimato per media campionaria su $N = 16$ realizzazioni
indipendenti per ciascun fotogramma, e $\mathcal{T}$ compone le seguenti
trasformazioni:
\begin{itemize}
    \item guadagno cromatico moltiplicativo per canale in $[0.6,\, 1.4]$;
    \item offset di luminosità additivo in $[-0.2,\, +0.2]$;
    \item rumore gaussiano additivo a media nulla e deviazione standard $0.05$;
    \item rotazione nel piano in $[-15^\circ,\, +15^\circ]$;
    \item riscalamento isotropo in $[0.4,\, 0.9]$;
    \item rapporto d'aspetto indipendente sui due assi in $[0.7,\, 1.3]$;
    \item traslazione entro il $10\%$ della dimensione della superficie.
\end{itemize}
La variazione indipendente dei fattori di scala sui due assi modella la
componente di beccheggio del velivolo, che comprime la proiezione del bersaglio
lungo una sola direzione. L'aggregazione sposta l'obiettivo dal caso migliore al
valore atteso sulla famiglia di condizioni modellata, ed è il presupposto che
rende l'artefatto un candidato alla realizzazione fisica e non un esemplare
valido nel solo dominio digitale.

L'invarianza geometrica non esaurisce il vincolo, poiché la catena di
riproduzione agisce anche sul contenuto in frequenza della superficie.
Un'ottimizzazione non vincolata converge verso pattern ad alta frequenza
spaziale, efficaci contro il modello e non riproducibili: retinatura di stampa,
trama del tessuto e campionamento del sensore agiscono come filtri passa-basso
in cascata e sopprimono selettivamente la banda su cui l'attacco si appoggia. Lo
stesso effetto è osservabile già in ottimizzazione, poiché l'interpolazione
bilineare delle trasformazioni geometriche attenua le componenti a frequenza più
elevata. Si introduce quindi un termine di variazione totale
\footcite{thys_2019},
\begin{equation}
    \mathcal{L}_{TV}(\mathrm{P}) =
    \frac{1}{C\,H\,(W-1)} \sum_{c,i,j} \left| p_{c,i,j+1} - p_{c,i,j} \right|
    + \frac{1}{C\,(H-1)\,W} \sum_{c,i,j} \left| p_{c,i+1,j} - p_{c,i,j} \right|
    \label{eq:tv}
\end{equation}
dove $p_{c,i,j}$ è il valore della patch $\mathrm{P}$ nel canale $c$ alla
posizione di indici $i$ e $j$, mentre $C$, $H$ e $W$ sono numero di canali,
altezza e larghezza della superficie. Si adotta la formulazione anisotropa in
norma $L_1$, mediata separatamente lungo le due direzioni: rispetto alla forma
isotropa è priva di radice quadrata e quindi di singolarità del gradiente nelle
regioni in cui la superficie è localmente costante. Il termine entra nella
funzione obiettivo con peso $\lambda_{TV}$, che regola il compromesso fra
aggressività e stampabilità: valori elevati spostano la soluzione verso pattern
a bassa frequenza spaziale, più robusti alla riproduzione e meno efficaci contro
il modello.

\textbf{Vincoli di piattaforma e budget di ottimizzazione.}
\label{sec:accumulo}
L'intera sperimentazione è condotta su una singola piattaforma Apple Silicon M4
con $16$\,GB di memoria unificata, condivisa fra unità di calcolo generale e
acceleratore. Ciascun fotogramma richiede sedici composizioni e altrettanti
passaggi in avanti attraverso il rilevatore: la dimensione del lotto fisico è
vincolata a uno, e l'aggregazione su più fotogrammi è ottenuta per accumulo del
gradiente su quattro iterazioni consecutive prima di ciascun aggiornamento, con
dimensione di lotto efficace pari a quattro a occupazione invariata. Il
campionamento su griglia delle trasformazioni geometriche è mantenuto in
precisione a virgola mobile singola per instabilità numerica del gradiente in
precisione ridotta sul backend adottato; il corrispondente operatore
all'indietro, privo di implementazione nativa, è eseguito su unità di calcolo
generale mediante ricaduta automatica, mentre l'inferenza resta
sull'acceleratore. Il numero di bersagli considerati per fotogramma è limitato a
tre, selezionati per area decrescente del \textit{bounding box}: il vincolo
contiene il costo di composizione sui fotogrammi densamente popolati e concentra
il segnale dove il rilevatore dispone di evidenza sufficiente.

L'accumulo introduce una distinzione fra due contatori. Le \emph{iterazioni}
corrispondono ai fotogrammi elaborati, gli \emph{aggiornamenti} agli interventi
effettivi dell'ottimizzatore sui parametri. Il programma di decadimento del
passo di apprendimento, di tipo coseno, deve essere espresso nella seconda
unità: espresso nella prima, il suo orizzonte risulta dilatato del fattore di
accumulo e il passo permane prossimo al valore iniziale per la quasi totalità
della traiettoria. Nell'implementazione adottata l'orizzonte è calcolato come
rapporto fra budget di iterazioni e fattore di accumulo, ossia sul numero di
aggiornamenti \emph{attesi}.

Il numero di aggiornamenti eseguiti risulta inferiore a quello atteso, perché un
fotogramma nel quale nessun bersaglio supera la soglia di rilevanza dimensionale
definita nella Sezione \ref{sec:loss} ha peso nullo ed è scartato prima
dell'aggiornamento. La frazione di iterazioni utili dipende dalla distribuzione
delle taglie propria dell'insieme di dati ed è stata misurata per strumentazione
diretta del ciclo di ottimizzazione. Il decadimento del passo non completa
quindi il proprio orizzonte entro la durata del run; entità dello scarto e
verifica sono riportate nel \hyperref[cap5]{Capitolo Quinto}, e la questione è ripresa fra le
limitazioni.

Due elementi del protocollo vanno dichiarati. Un criterio di arresto anticipato,
basato sulla media mobile dell'obiettivo su una finestra di cento aggiornamenti,
è configurato come salvaguardia contro la divergenza e non è intervenuto in
alcuno dei run riportati: la pazienza impostata è comparabile al numero
complessivo di aggiornamenti disponibili. La patch valutata nel \hyperref[cap5]{Capitolo Quinto} corrisponde all'ultimo aggiornamento eseguito e non a quello di
minima media mobile; le due differiscono in misura trascurabile, e la prima è la
formulazione corretta.

\section{Funzione obiettivo}
\label{sec:loss}

La funzione obiettivo si compone di due termini,
\begin{equation}
    \mathcal{L} = -\ln\left(1 - \bar{c} + \varepsilon\right)
    + \lambda_{TV} \, \mathcal{L}_{TV}(\mathrm{P})
    \label{eq:loss_totale}
\end{equation}
dove $\bar{c}$ è una stima aggregata della confidenza attribuita dal
rilevatore alla classe bersaglio nella regione del soggetto\footnote{Nel codice
di ottimizzazione la quantità aggregata è ottenuta applicando una
trasformazione logistica al punteggio di classe restituito dalla testa di
rilevamento, che in \textit{Ultralytics} è già normalizzato nell'intervallo
unitario. La composizione delle due trasformazioni restringe il codominio
effettivo di $\bar{c}$ a $[0.5,\,\sigma(1))$ e comprime la scala rispetto al
punteggio del rilevatore. Poiché la trasformazione è monotona crescente, la
direzione del gradiente e l'ordinamento fra configurazioni restano invariati; i
valori riportati nella Sezione~\ref{sec:risultato_divario} sono espressi anche
sulla scala del rilevatore, ottenuta per inversione.}, $\varepsilon =
10^{-6}$ è una costante che previene la divergenza del logaritmo,
$\mathcal{L}_{TV}$ è il termine di variazione totale dell'Equazione \ref{eq:tv}
e $\lambda_{TV} = 0.1$ è il suo peso. Il valore di $\lambda_{TV}$ è stato
fissato a valle dell'osservazione che valori inferiori di tre ordini di
grandezza producono superfici a contenuto in frequenza troppo elevato per
sopravvivere all'interpolazione bilineare delle trasformazioni geometriche.

\textbf{Dalla formulazione a soglia alla formulazione asintotica.}
\label{sec:loss_asintotica}
Una formulazione a soglia del tipo \textit{Hinge Loss}, che penalizza la sola
eccedenza della confidenza aggregata rispetto a una soglia di rilevamento, si
dimostra degenere come obiettivo differenziabile: la sua derivata è
identicamente nulla al di sotto della soglia, quindi l'aggiornamento dei
parametri cessa appena la confidenza scende sotto di essa, indipendentemente dal
fatto che la confidenza residua sia prossima alla soglia o trascurabile. Il
punto stazionario così raggiunto è quello di minimo margine, nel quale il
rilevamento è appena soppresso e una variazione delle condizioni di osservazione
lo ripristina.

La formulazione adottata impiega invece una barriera logaritmica sul complemento
della confidenza. Il termine $-\ln(1 - \bar{c} + \varepsilon)$ è definito e
strettamente crescente in $\bar{c}$ su tutto l'intervallo, e la sua derivata
rispetto a $\bar{c}$ vale $(1 - \bar{c} + \varepsilon)^{-1}$. Ne derivano due
proprietà. Il gradiente non si annulla in alcun punto del dominio:
l'ottimizzazione esercita pressione anche quando la confidenza è già sotto la
soglia di rilevamento, e la spinge asintoticamente verso lo zero, con un
margine rispetto alle variazioni delle condizioni di osservazione. La
penalizzazione cresce inoltre in modo più che proporzionale al crescere della
confidenza residua, quindi i casi in cui il rilevatore mantiene evidenza forte
pesano in modo sproporzionato e lo sforzo di ottimizzazione si concentra dove il
rilevamento è più solido: è la proprietà richiesta da una perturbazione
universale, il cui requisito è l'efficacia sul caso avverso. La forma analitica
è priva di limite superiore; nell'implementazione il codominio ristretto
richiamato in nota limita la penalizzazione a un valore finito, senza alterare
l'ordinamento fra livelli di confidenza residua.

L'assenza di una soglia esplicita nell'obiettivo comporta che essa non figuri
fra gli iperparametri dell'ottimizzazione. La soglia resta definita in fase di
valutazione, dove decide se una predizione sia un rilevamento, e non
condiziona la direzione del gradiente.

\textbf{Aggregazione della confidenza sulle celle salienti.}
\label{sec:topk}
Il rilevatore produce, per ciascun fotogramma, un numero di predizioni pari alla
somma delle celle delle tre mappe di attivazione a passi di campionamento $8$,
$16$ e $32$ pixel. La definizione di $\bar{c}$ richiede di stabilire su quali di
tali predizioni la confidenza vada aggregata.

Il primo criterio è geometrico: si considerano le sole celle il cui centro cade
all'interno del \textit{bounding box} del soggetto, ottenendo un insieme
$\mathcal{M}$ di indici. La restrizione è necessaria, perché l'avversario non
controlla il resto della scena e una penalizzazione della confidenza fuori dal
bersaglio ottimizzerebbe una perturbazione non realizzabile. Non è però
sufficiente: l'impronta di un bersaglio ben risolto comprende, sulle tre scale,
dell'ordine di duecento celle, delle quali una minoranza porta evidenza forte,
mentre le restanti si collocano nella periferia del \textit{bounding box}, dove
il segnale è debole e l'insieme delle celle attive varia in modo scorrelato fra
fotogrammi successivi. Una media estesa a tutto $\mathcal{M}$ diluisce il
segnale utile in un fondo variabile e produce un gradiente dominato dal rumore.

L'aggregazione è quindi ristretta alle $K$ celle a confidenza più elevata
all'interno di $\mathcal{M}$. Detto $s_j$ il punteggio della classe bersaglio
associato alla cella $j$-esima, e detto $\mathcal{K} \subseteq \mathcal{M}$
l'insieme degli indici dei $K$ punteggi maggiori,
\begin{equation}
    \bar{c} = \frac{\sum_{j \in \mathcal{K}} w_j \, s_j}{\sum_{j \in \mathcal{K}} w_j}
    \label{eq:topk}
\end{equation}
dove $w_j \in [0,1]$ è il peso di rilevanza dimensionale del bersaglio che ha
generato la cella $j$-esima. Detta $h$ l'altezza in pixel del \textit{bounding
box} di tale bersaglio, il peso vale
\begin{equation}
    w(h) =
    \begin{cases}
        0 & h < 60 \\[2pt]
        \dfrac{h - 60}{90} & 60 \leq h < 150 \\[6pt]
        1 & h \geq 150
    \end{cases}
    \label{eq:peso_tattico}
\end{equation}
La soglia inferiore esclude dall'obiettivo i bersagli troppo poco risolti perché
una perturbazione pari al $20\%$ della loro area possa alterare l'esito
dell'inferenza; la rampa attenua il contributo dei bersagli intermedi. Il peso
agisce sul solo calcolo della media: la selezione delle celle avviene per
confidenza grezza, in modo che rifletta la distribuzione effettiva dell'evidenza
del rilevatore, e la semantica probabilistica di $\bar{c}$ è preservata. Un
bersaglio a peso nullo resta composto sul canvas, per consistenza fisica della
scena, e non contribuisce al gradiente; un fotogramma nel quale tutti i bersagli
hanno peso nullo non produce aggiornamento.

La formulazione generalizza l'aggregazione per massimo adottata in letteratura
per l'attacco a rilevatori di persone \footcite{thys_2019}, e la estende ai
fotogrammi con più bersagli: un massimo puro concentrerebbe l'intero gradiente
su un solo soggetto per fotogramma, mentre $K > 1$ ripartisce la pressione sul
nucleo di celle salienti di ciascuno.

La scala di $K$ è vincolata da un'analisi del comportamento del rilevatore in
assenza di perturbazione. Per ciascun fotogramma positivo dell'insieme di
validazione le celle sono ordinate per confidenza decrescente; mediante somma
cumulativa si ricostruiscono sensitività e specificità a livello di cella per
ogni valore di $K$, senza ripetere l'inferenza. Le due grandezze producono
criteri di ottimalità divergenti: la misura $F_1$ premia la copertura completa
dell'impronta geometrica e individua un optimum di ordine centinaia, mentre la
media geometrica $\sqrt{R_1 R_2}$ penalizza la diluizione nello sfondo e
individua un optimum di ordine decine. La divergenza quantifica l'estensione del
nucleo di celle a evidenza forte, che risulta un sottoinsieme ristretto
dell'impronta. Il valore $K = 20$ è compreso fra gli optimum misurati sui due
domini secondo il criterio della media geometrica, assunto come primario perché
bilancia le due metriche effettivamente impiegate dal simulatore. L'analisi
caratterizza il rilevatore non perturbato ed è una verifica di
plausibilità della scala di $K$, non un predittore dell'efficacia dell'attacco:
la relazione fra le due grandezze non è stata stabilita in questo lavoro.

\textbf{Obiettivo di ottimizzazione e metrica di valutazione.}
\label{sec:divario}
L'obiettivo dell'Equazione \ref{eq:loss_totale} e le metriche definite nella
Sezione \ref{sec:metriche} misurano proprietà distinte dello stesso sistema, e
la distinzione condiziona l'interpretazione dei risultati.

La quantità $\bar{c}$ è una media su $K$ celle di una griglia interna al
rilevatore, calcolata prima della soppressione dei non-massimi, mediata su
sedici realizzazioni della distribuzione di trasformazioni e valutata
sull'insieme di ottimizzazione. L'\textit{Evasion Rate} è un tasso di eventi
binari definito a livello di fotogramma, valutato sull'insieme di validazione
dopo la soppressione dei non-massimi e dopo l'applicazione di una soglia di
confidenza. Fra le due si interpongono tre trasformazioni non lineari.

La prima è l'aggregazione per massimo operata dalla soppressione dei
non-massimi: un rilevamento sopravvive se almeno una cella supera la soglia,
quindi la quantità che determina l'esito è $\max_j s_j$, mentre l'obiettivo
agisce sulla media delle $K$ maggiori. Le due coincidono per $K = 1$; per $K >
1$ la media è un surrogato conservativo del massimo, e una sua riduzione non si
traduce linearmente in una riduzione del massimo. La seconda è la soglia di
confidenza, che è una funzione a gradino: per una popolazione di predizioni
addensata in prossimità della soglia, una traslazione di piccola entità produce
una variazione del tasso di attraversamento di entità maggiore, con ampiezza
dipendente dalla densità della distribuzione nell'intorno della soglia. La terza
è la differenza di condizione di osservazione: $\bar{c}$ è un valore atteso su
una distribuzione che comprende riscalamenti fino a $0.4$ e rotazioni fino a
$15^\circ$, mentre la valutazione avviene su una singola realizzazione,
geometricamente più favorevole all'avversario.

Una flessione marginale di $\bar{c}$ si traduce pertanto in un incremento
drastico dell'\textit{Evasion Rate}, senza che ciò costituisca un'incoerenza
matematica. L'efficacia operativa dell'attacco non deriva da un collasso
assoluto della confidenza, bensì dall'attraversamento al ribasso della soglia
decisionale da parte di una popolazione di predizioni già in condizioni di
criticità.

\section{Metriche di valutazione}
\label{sec:metriche}

L'unità di analisi è il fotogramma. Un bersaglio è valido quando l'altezza del
suo \textit{bounding box} annotato è pari o superiore a $60$ pixel nello spazio
del canvas di rete, soglia che esclude i bersagli per i quali il campionamento
spaziale comprometterebbe l'integrità della perturbazione. Un fotogramma è
positivo quando contiene almeno un bersaglio valido, negativo altrimenti.

\textbf{Classificazione degli esiti a livello di fotogramma.}
\label{sec:esiti}
Sui fotogrammi positivi l'esito è determinato dalla presenza di almeno un
rilevamento della classe bersaglio sopra soglia di confidenza: la presenza
produce un vero positivo, l'assenza un falso negativo. Nessun criterio di
sovrapposizione geometrica interviene, poiché l'obiettivo è stabilire se il
sistema rilevi la presenza di una persona nella scena, prescindendo
dall'accuratezza della localizzazione.

Sui fotogrammi negativi la classificazione impiega una convenzione a regione
ignorata. Un fotogramma negativo può contenere persone annotate di altezza
inferiore alla soglia di validità: un rilevamento su una di esse è corretto, ma
riferito a un bersaglio non tatticamente rilevante, e non è un errore.
Ciascun rilevamento è quindi confrontato per intersezione su unione con tutte le
persone annotate del fotogramma, di qualunque dimensione; un rilevamento con
sovrapposizione pari o superiore a $0.3$ è ignorato. Il fotogramma è classificato
come falso positivo se almeno un rilevamento resta privo di corrispondenza, come
vero negativo altrimenti.

La soglia di sovrapposizione adottata è più permissiva dello standard $0.5$
della campagna PASCAL VOC \footcite{everingham_2010_voc}. La scelta è motivata
dalla sensibilità della metrica su riquadri di piccole dimensioni, per i quali
uno scostamento di pochi pixel produce una variazione sproporzionata del valore
\footcite{yu_2020_scalematch}, ed è allineata alla convenzione impiegata su
insiemi di dati con occlusione elevata \footcite{shao_2018_crowdhuman}. La
convenzione è strettamente necessaria: l'esclusione asimmetrica dei fotogrammi
contenenti unicamente persone sotto soglia ridurrebbe la classe negativa a una
frazione marginale dell'insieme di validazione, e renderebbe degenere la stima
della specificità.

\textbf{Metriche primarie e secondarie.}
\label{sec:metriche_scelta}
Le metriche primarie sono la sensitività $R_1$ e la specificità $R_2$ definite
nell'Equazione \ref{eq:r1_r2}, l'\textit{Evasion Rate} e la loro media
geometrica. L'\textit{Evasion Rate} è il complemento a uno della sensitività,
$E = \mathit{FN} / (\mathit{TP} + \mathit{FN}) = 1 - R_1$, ed esprime la frazione
di fotogrammi contenenti un bersaglio valido nei quali il rilevatore non produce
alcun rilevamento: è la quantità che misura direttamente l'efficacia della
perturbazione. La media geometrica $G = \sqrt{R_1 R_2}$ aggrega le due
sensitività attribuendo peso identico alle due classi e si annulla se una delle
due si annulla. È assunta come indicatore primario aggregato perché garantisce il
bilanciamento fra classe positiva e negativa e si allinea esattamente ai due
parametri probabilistici acquisiti dal simulatore attraverso il file di
interscambio.

La misura $F_1$ è riportata come metrica secondaria. La sua precisione dipende
dal numero di falsi positivi, che sulla classe negativa estesa è governato dalla
convenzione a regione ignorata e non dall'effetto della perturbazione: $F_1$
risulta pertanto meno interpretabile delle due sensitività separate, e la sua
inadeguatezza come parametro del canale visivo del simulatore è già stata
argomentata nella Sezione \ref{sec:laws_sim}.

\textbf{Metriche accessorie.}
\label{sec:metriche_accessorie}
Il conteggio di rilevamenti spuri quantifica le rilevazioni fantasma indotte dalla
perturbazione: sui soli fotogrammi positivi, dove la patch è effettivamente
presente nella scena, si contano i rilevamenti privi di corrispondenza per
intersezione su unione con qualunque persona annotata. La quantità è un
conteggio grezzo per fotogramma e non un indicatore binario, poiché un
fotogramma con tre allucinazioni e uno con una sola non sono equivalenti.

Una seconda coppia di conteggi è calcolata nello stesso ciclo di inferenza,
ristretta ai bersagli di altezza pari o superiore a $80$ pixel, corrispondente
allo scenario di ingaggio ravvicinato modellato dal simulatore. La copertura
tattica è definita come frazione delle annotazioni valide che soddisfano anche
tale soglia, e non è confrontabile fra insiemi di dati valutati a risoluzioni
diverse.

Il simulatore calcola infine una metrica di ingaggio a costo pesato, che
rapporta gli ingaggi corretti agli errori di entrambi i tipi ponderati per il
costo dell'attacco che li produce. La sua formulazione e i valori che ne
derivano non sono riportati fra i risultati, per le ragioni esposte nella
Sezione~\ref{sec:esperimenti_sim}, e la metrica è ripresa fra gli sviluppi
futuri.

\section{Insiemi di dati}
\label{sec:dataset}

La valutazione impiega due insiemi di dati aerei, trattati separatamente e non
aggregati. L'impiego di due domini indipendenti è parte del disegno
sperimentale: la replica dell'effetto su distribuzioni di scala e condizioni di
ripresa distinte discrimina fra una vulnerabilità caratteristica di un singolo
insieme di dati e una proprietà del rilevatore. In entrambi i casi
ottimizzazione e valutazione impiegano suddivisioni disgiunte.

\textbf{VisDrone2019-DET.}
\label{sec:visdrone}
L'insieme di validazione comprende $531$ immagini statiche acquisite da velivolo
senza pilota in contesto urbano \footcite{zhu_2021_visdrone}, ridimensionate a
$640 \times 640$ pixel. La composizione misurata alla soglia di validità di $60$
pixel è la seguente: $80$ fotogrammi positivi, con almeno un bersaglio valido,
pari al $15.1\%$; $9$ fotogrammi privi di qualunque persona annotata, pari
all'$1.7\%$; $442$ fotogrammi contenenti soltanto persone sotto soglia, pari
all'$83.2\%$. La terza categoria motiva la convenzione a regione ignorata della
Sezione \ref{sec:metriche}: senza di essa la classe negativa si ridurrebbe ai
soli nove fotogrammi della seconda. Il numero medio di bersagli validi per
fotogramma è $0.72$.

\textbf{Okutama-Action.}
\label{sec:okutama}
Il secondo insieme comprende sequenze video acquisite da velivolo senza pilota
ad altitudine compresa fra $10$ e $45$ metri \footcite{barekatain_2017_okutama},
intervallo coerente con lo scenario di ingaggio ravvicinato modellato dal
simulatore, mentre la prospettiva di VisDrone corrisponde a una sorveglianza ad
area più estesa. L'insieme di validazione presenta il $57.0\%$ di fotogrammi
positivi e un numero medio di bersagli validi per fotogramma pari a $1.45$.

Sono impiegati i fotogrammi pre-estratti a $1280 \times 720$ pixel e le
annotazioni ad azione singola. Le coordinate delle annotazioni sono espresse
nello spazio nativo $3840 \times 2160$ e non nello spazio dei fotogrammi
distribuiti: il riscalamento verso il canvas di rete avviene in un unico
passaggio, senza stadio intermedio. Le colonne relative alla classificazione
dell'azione sono ignorate, secondo l'indicazione della documentazione ufficiale
per il compito di rilevamento di pedoni. La suddivisione impiegata è quella
ufficiale, che separa le sequenze per scenario e garantisce quindi che nessuna
sequenza compaia sia in ottimizzazione sia in valutazione.

\textbf{Risoluzione del canvas e vincolo di memoria.}
\label{sec:risoluzione}
La risoluzione di $960 \times 960$ pixel adottata su Okutama-Action deriva da un
vincolo verificato sperimentalmente. A $1280 \times 1280$ pixel, con sedici
trasformazioni per fotogramma e precisione a virgola mobile singola,
l'occupazione di memoria residente eccede la capacità della piattaforma: il
processo entra in condizione di \textit{Out-Of-Memory}, con paginazione
continua, e non completa l'ottimizzazione. A $960 \times 960$ l'occupazione
rientra nel budget, il contatore di paginazione resta nullo e il budget di
iterazioni è portato a termine.

La parità esatta con la risoluzione impiegata su VisDrone è stata scartata sulla
base di un calcolo sulla distribuzione delle taglie: a $640$ pixel ogni altezza
si dimezza rispetto a $1280$, quindi un riquadro resta valido solo se misurava
almeno $120$ pixel alla risoluzione superiore. Tale condizione elimina l'intero
intervallo compreso fra $60$ e $100$ pixel, che è la quasi totalità
della distribuzione misurata, riduce il numero di riquadri validi di oltre il
novanta percento e ripristina il problema che la migrazione al secondo insieme
di dati intendeva rimuovere. Il costo della risoluzione adottata è una riduzione
del numero di riquadri validi rispetto a $1280$ pixel, riportata fra le
limitazioni del lavoro.

\section{Protocollo statistico}
\label{sec:protocollo}

La significatività delle differenze fra condizione non perturbata e condizione
sotto attacco è stabilita per ricampionamento, senza assunzioni distribuzionali
sulle metriche.

\textbf{Intervalli di confidenza per ricampionamento.}
\label{sec:bootstrap}
L'unità di ricampionamento è il fotogramma. Per ciascun fotogramma la
valutazione produce un vettore a componente unica fra vero positivo, falso
negativo, falso positivo e vero negativo, secondo la classificazione della
Sezione \ref{sec:metriche}. L'insieme di tali vettori è ricampionato con
reinserimento per $10\,000$ iterazioni; a ogni iterazione i conteggi sono
riaggregati e la metrica ricalcolata. L'intervallo di confidenza al $95\%$ è
dato dal $2.5$-esimo e dal $97.5$-esimo percentile della distribuzione ottenuta,
secondo il metodo dei percentili \footcite{efron_1993_bootstrap}. Il valore
puntuale riportato è quello calcolato sul campione originale e non la media
delle repliche.

Il confronto fra le due condizioni impiega un ricampionamento appaiato: a ogni
iterazione una sola lista di indici è estratta e impiegata per riaggregare i
conteggi di entrambe le condizioni, che si riferiscono allo stesso insieme di
fotogrammi nello stesso ordine. La distribuzione ottenuta è quella della
differenza fra le metriche, e da essa si derivano l'intervallo di confidenza
della differenza e un valore di probabilità a due code sotto l'ipotesi nulla di
differenza nulla, calcolato come doppio della frazione minima di repliche che
contraddicono il segno osservato. La differenza è dichiarata significativa
quando lo zero non è compreso nell'intervallo. Il ricampionamento appaiato
rimuove la componente di variabilità comune alle due condizioni, che un
confronto fra intervalli calcolati indipendentemente non separa.

\textbf{Decorrelazione temporale.}
\label{sec:decorrelazione}
Il ricampionamento assume unità indipendenti. L'assunzione è soddisfatta dal
primo insieme di dati, costituito da immagini statiche distinte, e non dal
secondo, costituito da sequenze video acquisite a circa trenta fotogrammi al
secondo, nel quale fotogrammi consecutivi differiscono di pochi pixel e non
sono osservazioni distinte. L'impiego dell'insieme completo produce
pseudo-replicazione: le stime puntuali restano valide, l'incertezza è
sottostimata e gli intervalli risultano artificialmente stretti.

La correzione adottata è un sottocampionamento temporale con passo costante,
applicato all'indice del loader prima della valutazione. Il passo impiegato è di
$27$ fotogrammi, che riduce l'insieme di validazione a $527$ fotogrammi separati
da circa nove decimi di secondo e numericamente comparabili ai $531$ fotogrammi
del primo insieme. L'effetto della correzione è quantificato nel \hyperref[cap5]{Capitolo Quinto}: le stime puntuali coincidono al secondo decimale,
l'ampiezza degli intervalli aumenta di un fattore prossimo a cinque, e le
conclusioni di significatività non cambiano. I valori riportati come definitivi
sono quelli calcolati sull'insieme decorrelato.

\section{Disegno degli esperimenti}
\label{sec:esperimenti}

Gli esperimenti sono organizzati su tre piani: la validazione del canale visivo
su dati reali, l'analisi della dipendenza dell'efficacia dalla scala del
bersaglio, e la propagazione del degrado misurato attraverso la catena
decisionale simulata. Solo il primo produce misure sperimentali nel senso
definito nella Sezione \ref{sec:laws_sim}.

\textbf{Campagne di ottimizzazione e validazione.}
\label{sec:campagne}
Ciascuna campagna consiste nell'ottimizzazione di una patch universale sulla
suddivisione di addestramento di un insieme di dati, seguita dalla valutazione
sulla suddivisione di validazione del medesimo insieme. La patch ottimizzata su
un dominio non è impiegata per la valutazione sull'altro: le due campagne sono
indipendenti e i loro esiti non sono aggregati.

La valutazione confronta due condizioni sullo stesso insieme di fotogrammi e nel
medesimo ordine di iterazione. La condizione di controllo esegue l'inferenza sul
fotogramma non perturbato; la condizione sotto attacco compone la patch sulla
regione derivata dal \textit{bounding box} di ciascun bersaglio valido, secondo
la geometria definita nella Sezione \ref{sec:attacco}, e ripete l'inferenza.
L'appaiamento fra le due passate è la condizione che rende applicabile il
ricampionamento appaiato della Sezione \ref{sec:protocollo}. Sulla suddivisione
di validazione di Okutama-Action la valutazione è preceduta dal
sottocampionamento temporale con passo $27$, che agisce sull'indice del
\textit{loader} prima di entrambe le passate.

La soglia di confidenza del rilevatore è posta a $0.50$ per le misure riportate.
Le metriche accessorie definite nella Sezione \ref{sec:metriche} sono
ricalcolate anche alle soglie $0.30$ e $0.70$, per quantificare la dipendenza
dell'esito dal valore adottato e distinguere un effetto stabile da uno prodotto
dalla scelta di un particolare punto di taglio.

\textbf{Analisi della dipendenza dalla scala.}
\label{sec:stratificazione}
L'analisi verifica se l'efficacia della perturbazione dipenda dalla dimensione
apparente del bersaglio. L'unità di analisi è il singolo bersaglio e non il
fotogramma, e la stratificazione impiega tre intervalli di altezza del
\textit{bounding box} nello spazio del canvas: $[60, 100)$, $[100, 150)$ e
$[150, \infty)$ pixel. Il primo estremo coincide con la soglia di validità
definita nella Sezione \ref{sec:metriche}; il terzo intervallo raccoglie i bersagli meglio risolti, oltre la
soglia di copertura tattica definita nella Sezione~\ref{sec:metriche}.

L'esito è attribuito per corrispondenza geometrica: ciascun rilevamento della
classe bersaglio sopra soglia di confidenza è confrontato per intersezione su
unione con il \textit{bounding box} annotato del bersaglio in esame, e il
bersaglio è considerato rilevato quando almeno un rilevamento supera la soglia
di sovrapposizione $0.3$. L'attribuzione per corrispondenza è necessaria su
insiemi di dati con più soggetti simultanei nello stesso fotogramma, poiché un
criterio a livello di fotogramma classificherebbe come rilevato un bersaglio
effettivamente eluso ogni volta che il rilevatore individua un soggetto diverso
nella medesima scena, con una distorsione crescente al crescere della densità di
soggetti.

L'analisi è condotta in entrambe le condizioni, di controllo e sotto attacco,
sul medesimo insieme di bersagli. La quantità interpretabile è la differenza fra
le due, poiché il tasso misurato nella sola condizione sotto attacco non isola
l'effetto della perturbazione dal limite di localizzazione proprio del
rilevatore su bersagli di piccola dimensione apparente.

\textbf{Esperimenti di simulazione.}
\label{sec:esperimenti_sim}
Il simulatore opera su una griglia di $100$ metri di lato popolata da tre
bersagli e trenta entità civili, con velivolo ad altitudine di $10$ metri e
raggio operativo del sensore di $150$ metri. Ciascuno scenario è eseguito per
$150$ passi temporali. A ogni passo, per ciascuna entità entro il raggio
operativo, il canale visivo produce un esito secondo il modello di Bernoulli
della Sezione \ref{sec:laws_sim}, i tre canali sono combinati e la decisione è
classificata.

Sono definiti quattro scenari, per combinazione di due vettori di attacco
indipendenti: assenza di attacco; contromisura visiva attiva sui soli bersagli;
contaminazione dei profili informativi; entrambi gli attacchi simultanei. La
contaminazione agisce in due direzioni opposte: riduce gli indicatori di
rischio dei bersagli e innalza quelli delle entità civili, su tutti i campi
del profilo.

Due esperimenti impiegano gli stessi scenari con configurazioni di fusione
differenti. Il primo isola il canale visivo azzerando i pesi dei canali non
visivi, e quantifica l'effetto della sola contromisura fisica sulla decisione.
Il secondo impiega la ripartizione nominale, con peso $0.45$ al canale visivo,
$0.35$ al canale informativo e $0.20$ al canale comportamentale, e quantifica in
che misura la ridondanza fra canali attenui il degrado propagato dal primo
livello.

La classificazione degli esiti aggrega le azioni di ingaggio e di allerta in
un'unica categoria di intervento, poiché entrambe comportano l'attivazione della
catena operativa; le azioni di tracciamento passivo e di non intervento
costituiscono la categoria complementare. Il vincolo precauzionale descritto
nella Sezione \ref{sec:laws_sim} agisce a valle della classificazione per
soglia, e declassa l'ingaggio ad allerta quando l'ampiezza dell'intervallo di
confidenza mobile eccede $0.40$ oppure quando più di tre entità civili si
trovano in prossimità di un bersaglio con punteggio inferiore a $0.85$.

Ciascuno scenario è eseguito una sola volta, con generatore pseudocasuale
inizializzato a un valore fisso. Le quantità riportate per questi esperimenti
sono pertanto stime puntuali deterministiche e non dispongono di una misura di
incertezza campionaria, a differenza di quelle del primo livello. La limitazione
discende dallo statuto delle entità simulate e circoscrive l'uso di tali
quantità all'illustrazione del meccanismo di propagazione.

\chapter{Risultati e Discussione}
\label{cap5}

Il presente capitolo riporta i risultati del framework secondo l'asimmetria dichiarata nella Sezione~\ref{sec:laws_sim}: il canale visivo produce misure sperimentali, il simulatore multi-agente produce conseguenze di un modello. Le Sezioni~\ref{sec:risultati_ottimizzazione} e~\ref{sec:risultati_efficacia} riportano pertanto i soli risultati del Livello 1, la prima relativa al comportamento dell'ottimizzazione e la seconda all'efficacia della perturbazione risultante. La Sezione~\ref{sec:discussione} interpreta tali misure alla luce della letteratura e circoscrive lo statuto delle uscite del Livello 2/3.

\section{Ottimizzazione della perturbazione}
\label{sec:risultati_ottimizzazione}

\textbf{Budget di ottimizzazione e comportamento dell'ottimizzatore.}
\label{sec:budget_effettivo}
La frazione di iterazioni utili annunciata nella Sezione~\ref{sec:accumulo} è stata quantificata sul run di produzione su Okutama ($8000$ iterazioni, budget atteso di $2000$ aggiornamenti). Le proprietà misurate sono riportate in Tabella~\ref{tab:budget_okutama}.

\begin{table}[t]
\centering
\begin{tabular}{lr}
\toprule
Grandezza & Valore \\
\midrule
Aggiornamenti attesi / eseguiti & $2000$ / $1110$ ($55.5\%$) \\
Learning rate finale osservato / atteso ($\eta_{min}$) & $0.004726$ / $0.001$ \\
Media mobile della loss, minimo (aggiornamento $840$) & $0.7574$ \\
Media mobile della loss, finale & $0.7583$ \\
Norma del gradiente, intervallo & $6.8\times10^{-4}$ -- $3.0\times10^{-2}$ \\
Soglia di troncamento del gradiente (mai raggiunta) & $1.0$ \\
Scarto patch ultimo aggiornamento / miglior checkpoint & $0.0009$ \\
\bottomrule
\end{tabular}
\caption{Proprietà del run di ottimizzazione su Okutama-Action. Lo scarto fra aggiornamenti attesi ed eseguiti discende dal filtro di rilevanza tattica descritto nella Sezione~\ref{sec:loss}; la differenza fra patch finale e miglior checkpoint è espressa nella stessa unità della media mobile della loss.}
\label{tab:budget_okutama}
\end{table}

La strumentazione diretta del ciclo di ottimizzazione, condotta su un campione di $200$ iterazioni con contatori su tutti i punti di uscita, attribuisce l'intero scarto al filtro di rilevanza tattica ($75$ iterazioni su $200$, dello stesso ordine della frazione misurata sul run completo).

Il valore atteso dal programma coseno con $T_{max}=2000$ e $t=1110$ è $0.0047263$, coincidente alla sesta cifra decimale con il valore registrato: lo scarto rispetto a $\eta_{min}$ discende dal numero di aggiornamenti mancati. La media mobile della loss risulta piatta dall'aggiornamento $840$ al termine del run, e la soglia di troncamento del gradiente non è mai stata raggiunta: il criterio di arresto anticipato descritto nella Sezione~\ref{sec:accumulo} non interviene per costruzione. La patch valutata nel seguito del capitolo corrisponde all'ultimo aggiornamento eseguito; lo scarto rispetto al checkpoint di minima media mobile è trascurabile rispetto all'intervallo di variazione della loss durante il run.

\textbf{Validazione a posteriori della scala di aggregazione.}
\label{sec:risultati_topk}
Le Figure~\ref{fig:k_selection_visdrone} e~\ref{fig:k_selection_okutama} riportano le quattro grandezze definite nella Sezione~\ref{sec:loss} in funzione di $K$, calcolate senza patch sui rispettivi insiemi di validazione. Gli ottimi misurati sono riportati in Tabella~\ref{tab:k_ottimo}.

\begin{figure}[htbp]
\centering
\includegraphics[width=0.95\textwidth]{img/k_selection_plots_visdrone.pdf}
\caption[Selezione della scala di aggregazione su VisDrone]{Selezione di $K$ su VisDrone ($640 \times 640$): $F_1$, $\sqrt{R_1 R_2}$, $R_1$ e $R_2$ in funzione di $K$, insieme di validazione completo, nessuna perturbazione applicata. La linea tratteggiata indica l'ottimo secondo il criterio primario ($\sqrt{R_1 R_2}$, $K = 37$).}
\label{fig:k_selection_visdrone}
\end{figure}

\begin{figure}[htbp]
\centering
\includegraphics[width=0.95\textwidth]{img/k_selection_plots_okutama.pdf}
\caption[Selezione della scala di aggregazione su Okutama-Action]{Selezione di $K$ su Okutama-Action ($960 \times 960$), stessa procedura. La linea tratteggiata indica l'ottimo secondo il criterio primario ($K \approx 10$).}
\label{fig:k_selection_okutama}
\end{figure}

\begin{table}[t]
\centering
\begin{tabular}{lrr}
\toprule
Criterio & VisDrone ($640$) & Okutama ($960$) \\
\midrule
$F_1$ & $244$ & $\sim 55$ \\
$\sqrt{R_1 R_2}$ & $37$ & $\sim 10$ \\
\bottomrule
\end{tabular}
\caption{Valore di $K$ che massimizza ciascun criterio, per insieme di dati.}
\label{tab:k_ottimo}
\end{table}

Il valore adottato, $K = 20$, è compreso fra i due ottimi misurati secondo il criterio primario sui due domini. Il pannello relativo a $R_2$ di entrambe le figure mostra il crollo della specificità a livello di cella al crescere di $K$. La divergenza fra i due criteri misura l'estensione del nucleo di celle a evidenza forte rispetto all'impronta geometrica completa del bersaglio: l'ottimo secondo $F_1$ è di ordine centinaia su VisDrone e di ordine decine su Okutama, mentre l'ottimo secondo la media geometrica resta di ordine decine su entrambi.

\textbf{Convergenza delle sei configurazioni.}
\label{sec:sei_config}
Sei configurazioni dell'ottimizzazione, isolate una variabile alla volta, sono state confrontate sull'insieme di validazione VisDrone (Tabella~\ref{tab:sei_config}, Figura~\ref{fig:runs_comparison}).

\begin{table}[t]
\centering
\begin{tabular}{lrrrr}
\toprule
Configurazione & Aggiornamenti & Scheduler & $F_1$ & Evasion Rate \\
\midrule
0 (Hinge, storica) & $3000$ & corretto & $0.760$ & $38.7\%$ \\
1 (asintotica, media, acc.\ $=4$) & $375$ & allungato $4\times$ & $0.740$ & $41.25\%$ \\
2 (asintotica, media, acc.\ $=2$) & $750$ & allungato $2\times$ & $0.750$ & $40.0\%$ \\
3 (asintotica, media, acc.\ $=4$) & $2500$ & corretto & $0.720$ & $43.75\%$ \\
4 (asintotica, media, acc.\ $=16$) & $625$ & corretto & $0.740$ & $41.25\%$ \\
5 (asintotica, top-$K$, acc.\ $=4$) & $1250$ & corretto & $0.720$ & $43.75\%$ \\
\bottomrule
\end{tabular}
\caption{Sei configurazioni su VisDrone. La configurazione 5 raggiunge il massimo della configurazione 3 con metà del budget di aggiornamenti.}
\label{tab:sei_config}
\end{table}

\begin{figure}[htbp]
\centering
\includegraphics[width=0.95\textwidth]{img/runs_comparison.pdf}
\caption[Confronto delle sei configurazioni di ottimizzazione]{$F_1$ ed Evasion Rate per le sei configurazioni di Tabella~\ref{tab:sei_config}. Le etichette dell'asse orizzontale sono identificatori interni allo script di generazione e non la numerazione delle configurazioni impiegata nel testo; la corrispondenza segue lo stesso ordine di Tabella~\ref{tab:sei_config}.}
\label{fig:runs_comparison}
\end{figure}

Le sei configurazioni convergono in una banda compresa fra $38.7\%$ e $43.75\%$ di Evasion Rate. La norma del gradiente, misurata per fotogramma, risulta compresa fra $0.0007$ e $0.005$ in tutte e sei le configurazioni, e la soglia di troncamento non è mai raggiunta in alcun run, compreso quello su Okutama riportato in Tabella~\ref{tab:budget_okutama}. L'aumento del fattore di accumulo da $4$ a $16$, che distingue la configurazione 4 dalla configurazione 3, non produce un miglioramento del risultato.

\section{Efficacia dell'attacco avversariale}
\label{sec:risultati_efficacia}

\textbf{VisDrone2019-DET.}
\label{sec:risultato_visdrone}
Tabella~\ref{tab:risultato_visdrone} e Figura~\ref{fig:candele_visdrone} riportano le metriche primarie e la metrica secondaria su VisDrone ($n = 531$, soglia di confidenza $0.50$, ricampionamento appaiato su $10\,000$ iterazioni).

\begin{table}[t]
\centering
\begin{tabular}{lrrr}
\toprule
Metrica & PRE & POST & $\Delta$ [IC $95\%$] \\
\midrule
Evasion rate & $0.175$ & $0.425$ & $+0.250$ [$+0.156$, $+0.346$] \\
$R_1$ & $0.825$ & $0.575$ & $-0.250$ [$-0.346$, $-0.156$] \\
$R_2$ & $0.969$ & $0.969$ & $0.000$ [$0.000$, $0.000$] \\
$\sqrt{R_1 R_2}$ & $0.894$ & $0.746$ & $-0.148$ [$-0.210$, $-0.090$] \\
$F_1$ (secondaria) & $0.825$ & $0.657$ & $-0.168$ [$-0.242$, $-0.100$] \\
\bottomrule
\end{tabular}
\caption{Risultato consolidato su VisDrone, $n = 531$. Il valore di probabilità è inferiore a $0.0001$ per tutte le metriche a eccezione di $R_2$, per la quale vale $1.0$.}
\label{tab:risultato_visdrone}
\end{table}

\begin{figure}[htbp]
\centering
\includegraphics[width=0.95\textwidth]{img/fig_candele_visdrone_640.pdf}
\caption[Stime puntuali e differenza appaiata su VisDrone]{Stime puntuali nelle condizioni pre e post attacco con intervallo di confidenza al $95\%$ (a); differenza appaiata post meno pre con il relativo intervallo (b). Insieme di dati VisDrone, $n = 531$.}
\label{fig:candele_visdrone}
\end{figure}

\textbf{Robustezza alla soglia di confidenza.}
\label{sec:multisoglia}
Le metriche primarie sono state ricalcolate su VisDrone alle soglie $0.30$ e
$0.70$, oltre al valore di riferimento $0.50$. Evasion Rate, $R_1$,
$\sqrt{R_1 R_2}$ e $F_1$ restano significativi a tutte e tre le soglie, con
valori di probabilità compresi fra $0.0002$ e valori inferiori a $0.0001$. La
specificità $R_2$ risulta identica fra condizione pre e post a ciascuna soglia,
con valori rispettivamente di $0.825$, $0.969$ e $1.000$: l'invarianza
discussa più avanti non è quindi un artefatto del particolare punto di taglio
adottato.

La condizione di controllo mostra una dipendenza dalla soglia di entità
superiore all'effetto dell'attacco stesso. In assenza di qualunque
perturbazione, l'Evasion Rate su VisDrone passa dal $3.75\%$ alla soglia $0.30$
al $72.5\%$ alla soglia $0.70$. Il valore di riferimento $0.50$ adottato per
tutte le misure riportate si colloca in posizione intermedia rispetto a questo
intervallo.

\textbf{Okutama-Action.}
\label{sec:risultato_okutama}
Tabella~\ref{tab:risultato_okutama} e Figura~\ref{fig:candele_okutama} riportano le stesse metriche su Okutama, sull'insieme decorrelato ($n = 527$, Sezione~\ref{sec:protocollo}).

\begin{table}[t]
\centering
\begin{tabular}{lrrr}
\toprule
Metrica & PRE & POST & $\Delta$ [IC $95\%$] \\
\midrule
Evasion rate & $0.497$ & $0.777$ & $+0.280$ [$+0.228$, $+0.332$] \\
$R_1$ & $0.503$ & $0.223$ & $-0.280$ [$-0.332$, $-0.228$] \\
$R_2$ & $0.960$ & $0.960$ & $0.000$ [$0.000$, $0.000$] \\
$\sqrt{R_1 R_2}$ & $0.695$ & $0.463$ & $-0.232$ [$-0.278$, $-0.188$] \\
$F_1$ (secondaria) & $0.657$ & $0.356$ & $-0.300$ [$-0.358$, $-0.244$] \\
\bottomrule
\end{tabular}
\caption{Risultato consolidato su Okutama-Action, $n = 527$. Il valore di probabilità è inferiore a $0.0001$ per tutte le metriche a eccezione di $R_2$, per la quale vale $1.0$.}
\label{tab:risultato_okutama}
\end{table}

\begin{figure}[htbp]
\centering
\includegraphics[width=0.95\textwidth]{img/fig_candele_okutama_960_stride27.pdf}
\caption[Stime puntuali e differenza appaiata su Okutama-Action]{Stime puntuali nelle condizioni pre e post attacco con intervallo di confidenza al $95\%$ (a); differenza appaiata post meno pre con il relativo intervallo (b). Insieme di dati Okutama-Action, $n = 527$. La specificità $R_2$ è la sola metrica il cui intervallo della differenza collassa su un punto singolo nel pannello (b).}
\label{fig:candele_okutama}
\end{figure}

\textbf{Invarianza della specificità.}
\label{sec:r2_invariante}
In entrambi i domini $R_2$ risulta identico fra condizione pre e post fino alla sedicesima cifra decimale. La causa è strutturale: la patch è composta esclusivamente sui \textit{bounding box} dei bersagli validi (Sezione~\ref{sec:esperimenti}), quindi nei fotogrammi privi di bersaglio l'immagine sottoposta al rilevatore è identica nelle due condizioni, e l'inferenza di YOLOv8 è deterministica. La quantità misura il tasso di falso allarme di base del rilevatore sullo sfondo dell'insieme di dati, pari al $3.1\%$ su VisDrone e al $4.0\%$ su Okutama, e non un effetto della perturbazione. Ne discende che l'intero movimento di $\sqrt{R_1 R_2}$ è attribuibile a $R_1$. La replica su un secondo dominio indipendente esclude che l'invarianza sia una proprietà accidentale del primo insieme di dati.

\textbf{Effetto della decorrelazione temporale.}
\label{sec:risultato_decorrelazione}
Il confronto fra l'insieme correlato ($n = 14210$) e quello decorrelato ($n = 527$, Sezione~\ref{sec:protocollo}) è riportato in Tabella~\ref{tab:decorrelazione}.

\begin{table}[t]
\centering
\begin{tabular}{lrr}
\toprule
Metrica & $n = 14210$ (correlato) & $n = 527$ (passo $27$) \\
\midrule
Evasion rate PRE & $0.4988$ & $0.4967$ \\
Evasion rate POST & $0.7800$ & $0.7767$ \\
$\Delta$ Evasion rate [IC $95\%$] & $+0.2812$ [$+0.2712$, $+0.2909$] & $+0.2800$ [$+0.2281$, $+0.3322$] \\
\bottomrule
\end{tabular}
\caption{Effetto della decorrelazione temporale: stima puntuale sostanzialmente invariata, ampiezza dell'intervallo di confidenza circa cinque volte maggiore.}
\label{tab:decorrelazione}
\end{table}

Le stime puntuali coincidono al secondo decimale; l'ampiezza dell'intervallo aumenta di un fattore prossimo a cinque. La significatività non cambia, con valore di probabilità inferiore a $0.0001$ in entrambi i casi: la correzione modifica la stima dell'incertezza e non la conclusione.

\textbf{Rilevamenti spuri.}
\label{sec:risultato_collaterale}
Tabella~\ref{tab:collaterale} riporta il conteggio medio di rilevamenti privi di corrispondenza per fotogramma positivo (Sezione~\ref{sec:metriche}), alle tre soglie di confidenza.

\begin{table}[t]
\centering
\begin{tabular}{llrrrl}
\toprule
Insieme di dati & Soglia & PRE & POST & $\Delta$ & $p$ \\
\midrule
VisDrone & $0.30$ & $0.4625$ & $0.4375$ & $-0.025$ & $0.535$ \\
VisDrone & $0.50$ & $0.1125$ & $0.0500$ & $-0.063$ & $0.0136$ \\
VisDrone & $0.70$ & $0.0000$ & $0.0000$ & $0.000$ & $1.000$ \\
\addlinespace
Okutama & $0.30$ & $0.3700$ & $0.2867$ & $-0.083$ & $0.0108$ \\
Okutama & $0.50$ & $0.0967$ & $0.0433$ & $-0.053$ & $<0.0001$ \\
Okutama & $0.70$ & $0.0000$ & $0.0000$ & $0.000$ & $1.000$ \\
\bottomrule
\end{tabular}
\caption{Conteggio di rilevamenti spuri nelle condizioni pre e post attacco, per soglia di confidenza. Su VisDrone la differenza è significativa alla sola soglia $0.5$ e non sopravvive alla correzione di Bonferroni per sei confronti, che richiede $p < 0.0083$: il risultato è ritirato. Su Okutama la differenza sopravvive alla medesima correzione alla soglia $0.5$ ($p < 0.0001$); alla soglia $0.3$ è significativa al livello convenzionale $0.05$ ($p = 0.0108$) ma non alla soglia corretta; alla soglia $0.7$ non è misurabile per assenza di eventi nel campione decorrelato ($n = 300$), mentre una controprova sul campione completo correlato ($n = 8095$) conferma la stessa direzione con $p = 0.0038$.}
\label{tab:collaterale}
\end{table}

Il segno della differenza è negativo su entrambi i domini: la perturbazione riduce il numero di rilevamenti spuri. La Figura \ref{fig:prima_dopo} riporta due fotogrammi che illustrano il fenomeno: in entrambi la soppressione riguarda il solo soggetto perturbato, mentre le persone non perturbate presenti nella medesima scena restano rilevate.

\begin{figure}[htbp]
\centering
\includegraphics[width=0.85\textwidth]{img/prima_dopo_specificita.pdf}
\caption{Effetto della perturbazione su due fotogrammi degli insiemi di
validazione. In ciascuna coppia, a sinistra la condizione di controllo con il
riquadro del bersaglio e la relativa confidenza, a destra la condizione sotto
attacco. (a) Okutama-Action: la confidenza sul bersaglio passa da $0.73$ a
$0.00$ e le cinque persone non perturbate presenti nella scena restano rilevate
in entrambe le condizioni. (b) VisDrone: la confidenza passa da $0.66$ a $0.00$
con quattro persone non perturbate conservate. La soppressione riguarda il solo
soggetto su cui la perturbazione è composta, in accordo con il segno negativo
del conteggio di rilevamenti spuri riportato in
Tabella~\ref{tab:collaterale}. Fotogrammi da Okutama-Action
\footcite{barekatain_2017_okutama} e VisDrone \footcite{zhu_2021_visdrone},
entrambi distribuiti con licenza CC BY-NC-SA 3.0.}
\label{fig:prima_dopo}
\end{figure}

\textbf{Dipendenza dalla scala del bersaglio.}
\label{sec:risultato_stratificazione}
La misura con condizione di controllo e attribuzione per corrispondenza geometrica (Sezione~\ref{sec:esperimenti}) è riportata in Tabella~\ref{tab:stratificazione} per Okutama, unico insieme di dati sul quale l'analisi è stata condotta.

\begin{table}[t]
\centering
\begin{tabular}{lrrrr}
\toprule
Intervallo di altezza (a $960$ px) & Controllo & Attacco & $\Delta$ & $n$ \\
\midrule
$[60, 100)$ px & $70.0\%$ & $97.0\%$ & $+27.0$ pp & $19921$ \\
$[100, 150)$ px & $97.1\%$ & $100.0\%$ & $+2.9$ pp & $618$ \\
$[150, \infty)$ px & -- & -- & -- & $0$ \\
\bottomrule
\end{tabular}
\caption{Tasso di mancata localizzazione per bersaglio, per intervallo di altezza. Il terzo intervallo è vuoto per costruzione: a $960$ pixel una soglia di $150$ pixel equivale a $200$ pixel alla risoluzione di $1280$.}
\label{tab:stratificazione}
\end{table}

L'effetto è concentrato nell'intervallo $[60, 100)$ pixel, che contiene il $97\%$ dei bersagli validi. L'intervallo $[100, 150)$ pixel è in saturazione: la condizione di controllo è già al $97.1\%$ e non resta margine misurabile. Il delta a livello di bersaglio, pari a $+27.0$ punti percentuali, è in accordo con il delta a livello di fotogramma misurato indipendentemente e riportato in Tabella~\ref{tab:risultato_okutama}, pari a $+28.0$ punti percentuali: due unità di analisi e due criteri di attribuzione distinti convergono entro un punto percentuale.

\textbf{Divario fra obiettivo di ottimizzazione e metrica di valutazione.}
\label{sec:risultato_divario}
La quantità $\bar{c}$, ricavata per inversione dell'Equazione~\ref{eq:loss_totale}, varia da $0.5465$ a $0.5312$ su Okutama. Riportata sulla scala del punteggio del rilevatore, per inversione della trasformazione logistica richiamata nella Sezione~\ref{sec:loss}, la medesima variazione corrisponde a un passaggio da $0.187$ a $0.125$, pari a un delta di $-6.2$ punti percentuali. Il delta di Evasion Rate sullo stesso dominio è di $+28.0$ punti percentuali, con un rapporto approssimativo di $1{:}4.5$ fra le due variazioni. L'ordine di grandezza è quello atteso dalle tre trasformazioni non lineari formalizzate nella Sezione~\ref{sec:divario}. La medesima compressione di scala rende conto dell'intervallo entro cui si colloca la norma del gradiente riportata in Tabella~\ref{tab:budget_okutama}, e della circostanza che la soglia di troncamento non sia mai stata raggiunta in alcun run.

\textbf{Confronto fra i due domini.}
\label{sec:confronto_complessivo}
Tabella~\ref{tab:confronto_complessivo} accosta le misure ottenute sui due insiemi di dati, valutati separatamente e con patch distinte.

\begin{table}[t]
\centering
\small
\begin{tabular}{lrr}
\toprule
Metrica & VisDrone ($n = 531$, $640$ px) & Okutama ($n = 527$, $960$ px) \\
\midrule
Evasion rate PRE $\to$ POST & $0.175 \to 0.425$ & $0.497 \to 0.777$ \\
$\Delta$ Evasion rate [IC $95\%$] & $+0.250$ [$+0.156$, $+0.346$] & $+0.280$ [$+0.228$, $+0.332$] \\
$R_1$ PRE $\to$ POST & $0.825 \to 0.575$ & $0.503 \to 0.223$ \\
$R_2$ & $0.969$ (invariante) & $0.960$ (invariante) \\
$\Delta \sqrt{R_1 R_2}$ & $-0.148$ & $-0.232$ \\
$\Delta F_1$ & $-0.168$ & $-0.300$ \\
Rilevamenti spuri & ritirato & significativo a $0.5$ \\
Fotogrammi positivi & $15.1\%$ & $57.0\%$ \\
Bersagli validi per fotogramma & $0.72$ & $1.45$ \\
$K$ ottimo ($F_1$ / media geometrica) & $244$ / $37$ & $\sim 55$ / $\sim 10$ \\
\bottomrule
\end{tabular}
\caption{Confronto fra i due domini. La copertura tattica alla soglia di $80$ pixel è esclusa dal confronto: non è confrontabile fra insiemi di dati valutati a risoluzioni diverse, per le ragioni esposte nella Sezione~\ref{sec:metriche}.}
\label{tab:confronto_complessivo}
\end{table}

\section{Discussione dei risultati}
\label{sec:discussione}

Il contributo marginale della perturbazione è quasi identico sui due domini, pari a $+25$ e $+28$ punti percentuali di Evasion Rate, a fronte di linee di base e distribuzioni di scala distinte (Tabella~\ref{tab:confronto_complessivo}). Ciò che cambia fra i due insiemi di dati è la condizione di partenza del rilevatore, non l'entità dell'effetto attribuibile all'attacco. Le sei configurazioni della Sezione~\ref{sec:risultati_ottimizzazione} variano formulazione della loss, criterio di aggregazione, fattore di accumulo e programma di decadimento del passo, e convergono in una banda di cinque punti percentuali. La norma del gradiente misurata sulle sei configurazioni, compresa fra $0.0007$ e $0.005$, e il mancato intervento del troncamento in alcun run collocano il vincolo nel segnale disponibile e non nella capacità dell'ottimizzatore: un gradiente troppo elevato avrebbe attivato il troncamento, e un budget insufficiente sarebbe stato compensato dall'aumento del fattore di accumulo, che invece non produce guadagno.

La letteratura sugli attacchi avversariali in prospettiva aerea è coerente con questa lettura. I lavori che riportano tassi di successo elevati su VisDrone selezionano come bersaglio i veicoli \footcite{shrestha_2023}; l'ottimizzazione multi-dimensionale delle caratteristiche scarta esplicitamente la classe pedone per l'esiguità della sua estensione in pixel da vista zenitale \footcite{xi_2025_lfrap}; gli attacchi efficaci su immagini aeree operano su bersagli di dimensione maggiore e agiscono sulle rappresentazioni intermedie della rete anziché sull'uscita finale \footcite{tang_2023}. La formulazione per aggregazione delle celle a evidenza più elevata è ripresa dagli attacchi a rilevatori di persone a livello del suolo \footcite{thys_2019}, dove il soggetto occupa una frazione dell'immagine di due ordini di grandezza superiore. Il vincolo di copertura del $20\%$ dell'area del bersaglio, imposto dalla realizzabilità come capo di abbigliamento \footcite{huang_2020_upc}, agisce su una superficie che in prospettiva aerea misura poche decine di pixel per lato. L'andamento della media mobile della loss, piatta dall'aggiornamento $840$ su un budget di $1110$, è compatibile con i rendimenti marginali decrescenti osservati per le patch avversariali universali \footcite{brown_2017_patch}. Il limite osservato è quindi attribuibile alla dimensione in pixel del bersaglio in prospettiva aerea, e la sua replica su un secondo dominio indipendente lo qualifica come proprietà del compito.

La misura riportata in Tabella~\ref{tab:stratificazione} colloca l'effetto nell'intervallo dimensionale dominante e chiude una questione che su VisDrone era rimasta aperta, per l'esiguità dei casi disponibili nell'intervallo dimensionale superiore. L'accordo entro un punto percentuale fra il delta per bersaglio e il delta per fotogramma, ottenuti con unità di analisi e criteri di attribuzione distinti, è una verifica di coerenza interna fra due misure indipendenti. Il conteggio di rilevamenti spuri, ritirato su VisDrone per mancata robustezza multi-soglia, risulta significativo su Okutama alla soglia $0.5$, con un effetto della stessa direzione ma non significativo dopo correzione alla soglia $0.3$: è l'unica quantità per la quale il secondo dominio produce un risultato nuovo anziché una replica. Il suo segno negativo indica una perturbazione che sopprime il rilevamento del soggetto sul quale è applicata senza alterare il comportamento del rilevatore sul resto della scena.

Due elementi circoscrivono la portata del confronto. La quadruplicazione dei fotogrammi positivi, dal $15.1\%$ al $57.0\%$, non corrisponde a una potenza statistica quadruplicata: una volta rimossa la correlazione temporale (Tabella~\ref{tab:decorrelazione}) il campione effettivo di Okutama, pari a $527$ unità, è comparabile ai $531$ fotogrammi di VisDrone. Il guadagno del secondo insieme di dati è nel dominio osservato e non nella numerosità campionaria. La condizione di partenza su Okutama è a sua volta più debole, con un Evasion Rate del $49.7\%$ in assenza di attacco contro il $17.5\%$ di VisDrone. Fra le spiegazioni possibili figurano uno scarto di dominio del rilevatore rispetto alle posture riprese dall'alto e la distorsione introdotta dal ridimensionamento a canvas quadrato; nessuna delle due è stata verificata in questo lavoro. La verifica richiederebbe la stratificazione per taglia della sola condizione di controllo, non eseguita, ed è riportata fra gli sviluppi futuri.

Le uscite del simulatore multi-agente non sono presentate come risultato quantitativo. Quattro difetti di implementazione, accertati per lettura diretta del codice, ne condizionano l'interpretazione: l'esito del rilevamento non è un test statistico ma un'estrazione da distribuzione uniforme con soglia fissa; la condizione di riferimento non legge le metriche del Livello 1 in assenza di attacco, e confronta quindi una costante di letteratura con il dato misurato sotto attacco; la specificità non entra nella catena decisionale per le entità civili; l'effetto della distanza è conteggiato sia nella metrica empirica sia nel modello analitico, in contrasto con quanto stabilito nella Sezione~\ref{sec:laws_sim}. La correzione è definita nel suo perimetro e non è stata eseguita. L'unico enunciato sostenibile allo stato attuale è che il sistema a canali fusi mostra un degrado inferiore rispetto al canale visivo isolato, coerente con la funzione di ridondanza attribuita alla fusione; l'entità del fenomeno non è quantificabile con i dati disponibili.


\chapter{Conclusioni}
\label{cap6}

L'impiego di sistemi d'arma autonomi letali sposta la selezione del bersaglio da un operatore a un modello statistico, e con essa il presupposto di fatto su cui il Diritto Internazionale Umanitario fonda il principio di distinzione. Lo stallo dei negoziati internazionali, esaminato nel \hyperref[cap1]{Capitolo Primo}, e la dissoluzione della responsabilità individuale nella catena algoritmica, discussa nel \hyperref[cap2]{Capitolo Secondo}, lasciano il requisito di controllo umano significativo privo di un'effettiva verifica tecnica indipendente. Il presente lavoro ha affrontato questa lacuna per via sperimentale: ha sottoposto il canale sensoriale di un sistema di percezione a un attacco avversariale controllato e ne ha misurato il degrado con un protocollo statistico esplicito.

Il metodo è consistito nell'ottimizzazione di una perturbazione avversariale universale contro un rilevatore di oggetti della famiglia YOLO, applicata a una porzione del busto del soggetto compatibile con la realizzazione per mezzo di un capo di abbigliamento. La perturbazione è stata ottimizzata su una distribuzione di trasformazioni geometriche e fotometriche, così da conservare efficacia sotto condizioni di ripresa variabili anziché su una singola resa dell'immagine, e vincolata a una banda spaziale compatibile con la riproduzione a stampa. La valutazione ha impiegato due insiemi di dati aerei indipendenti, trattati separatamente e mai aggregati, e ha confrontato la medesima sequenza di fotogrammi nelle condizioni con e senza perturbazione. La significatività è stata stabilita per ricampionamento appaiato, con correzione della correlazione temporale sul secondo insieme di dati.

L'esito principale è un degrado statisticamente significativo della capacità del rilevatore di individuare la presenza di una persona, replicato su entrambi i domini con entità comparabile. Il tasso di falso allarme sugli sfondi privi di bersaglio è rimasto invece invariato, per una ragione strutturale del disegno sperimentale che è stata identificata e dichiarata: la perturbazione non è presente nei fotogrammi che quella quantità misura. L'effetto si concentra sui bersagli di piccola estensione apparente, che sono la quasi totalità della popolazione osservabile da quota aerea, mentre sui bersagli meglio risolti il rilevatore mostra un limite di localizzazione proprio, indipendente dall'attacco.

La lettura di questi risultati è la parte che eccede il singolo esperimento. Configurazioni di ottimizzazione radicalmente diverse per formulazione dell'obiettivo, criterio di aggregazione e disciplina del passo di apprendimento convergono su un esito quasi identico; la misura diretta del segnale di aggiornamento ne indica la debolezza intrinseca; la letteratura sugli attacchi in prospettiva aerea concorda nel selezionare bersagli di maggiore estensione. Il limite osservato non è quindi un difetto contingente di una particolare implementazione né una caratteristica di un singolo insieme di dati, ma una proprietà del compito: da quota, una figura umana occupa un numero di elementi di immagine tale da collocare il rilevamento in una condizione di criticità permanente. La conseguenza per il quadro normativo è diretta. Un sistema di targeting che opera su bersagli in questa condizione dispone di un margine di evidenza ridotto anche in assenza di ogni interferenza, e una perturbazione realizzabile con mezzi ordinari è sufficiente a spostarne l'esito. La verifica del principio di distinzione demandata a un canale percettivo con tali proprietà non offre le garanzie che la delega della forza letale presupporrebbe.

\section{Limitazioni}

Il regime di conoscenza adottato è massimamente favorevole all'attaccante: la perturbazione è ottimizzata sfruttando i gradienti di un modello surrogato assunto coincidente con il rilevatore bersaglio. Le misure riportate sono pertanto un limite superiore all'efficacia conseguibile e non una stima operativa; il disallineamento fra modello surrogato e sistema reale ne ridurrebbe l'entità in misura non quantificata.

L'attacco è stato valutato nel solo dominio digitale. La perturbazione è stata progettata con i vincoli richiesti dalla riproduzione fisica, ma la catena stampa-tessuto-illuminazione-sensore non è stata attraversata sperimentalmente, e nessuna delle misure riportate ne incorpora la degradazione.

La condizione di partenza del rilevatore sul secondo insieme di dati risulta sensibilmente più debole che sul primo, in assenza di qualunque attacco. La causa non è stata determinata: le ipotesi formulate riguardano lo scarto fra il dominio di addestramento del rilevatore e le posture riprese dall'alto, oppure la deformazione introdotta dal ridimensionamento dei fotogrammi, e nessuna delle due è stata verificata. Attaccare un rilevatore già compromesso è una condizione di prova meno pulita.

La correzione della correlazione temporale, necessaria su dati video, riduce il campione effettivo a una numerosità comparabile a quella del primo insieme di dati: il secondo dominio amplia la varietà delle condizioni osservate, non la potenza statistica disponibile. L'analisi della dipendenza dalla scala è stata condotta su un solo insieme di dati, e la fascia dei bersagli di maggiore estensione resta priva di osservazioni su entrambi. La risoluzione di lavoro è stata fissata da un vincolo di memoria della piattaforma, con una riduzione dichiarata del numero di bersagli utilizzabili.

I livelli di simulazione e di fusione dei metadati non producono risultati quantitativi. La loro esecuzione è deterministica e priva di misura di incertezza campionaria, e la lettura del codice ha accertato difetti di implementazione che ne condizionano l'interpretazione. Il loro contributo si limita all'illustrazione del meccanismo con cui un degrado percettivo si propaga lungo una catena decisionale ridondante.

\section{Sviluppi futuri: il C.A.R.E. Kit}

La prosecuzione naturale del lavoro è il passaggio al dominio fisico. La perturbazione caratterizzata in questa tesi è stata ottimizzata come valore atteso su una distribuzione di trasformazioni, che è precisamente la proprietà richiesta perché un artefatto stampato conservi efficacia sotto angolazioni, distanze e condizioni di luce variabili. La sua realizzazione come capo indossabile è la componente percettiva di un dispositivo più ampio, denominato C.A.R.E. Kit (\textit{Camouflage Against Recognition and Engagement}), concepito come strumento di autotutela per popolazioni civili esposte a sorveglianza e targeting automatizzati in contesti privi di regolamentazione effettiva. La linea si inserisce in un filone consolidato, che comprende il camuffamento asimmetrico contro il riconoscimento del volto \footcite{harvey_2010} e i tessuti a motivo avversariale \footcite{wu_2020, hu_2022}. Le implicazioni etiche della diffusione di contromisure di questo tipo, e la tensione fra tutela individuale ed elusione dei sistemi di sicurezza legittimi, sono state discusse in letteratura \footcite{helkala_2023} e vanno affrontate prima di qualunque proposta operativa.

Sul piano sperimentale restano aperte alcune verifiche a basso costo. La stratificazione per taglia della sola condizione di controllo discriminerebbe fra le due ipotesi formulate per la debolezza della linea di base sul secondo insieme di dati. L'estensione dell'obiettivo all'intera impronta geometrica del bersaglio, anziché al solo nucleo di celle a evidenza forte, richiede un nuovo ciclo di ottimizzazione ma è suggerita dalla divergenza fra i due criteri di selezione misurati. L'applicazione dell'obiettivo alle rappresentazioni intermedie della rete, anziché alla sua uscita, è la tecnica adottata dai lavori che riportano i risultati più solidi su immagini aeree. La correzione dei difetti accertati nei livelli di simulazione e fusione, infine, è definita nel suo perimetro e restituirebbe misura quantitativa alla propagazione del degrado lungo la catena decisionale, che allo stato attuale resta descritta e non misurata. \\
La stessa correzione consentirebbe inoltre di
riportare la metrica di ingaggio a costo pesato introdotta nel disegno
sperimentale, che rapporta gli ingaggi corretti agli errori di entrambi i tipi
ponderati per il costo dell'attacco che li produce. La sua taratura richiede
però di rendere commensurabili due costi oggi eterogenei, la frazione di
superficie corporea coperta dalla perturbazione e la frazione di campi
informativi contaminati, questione aperta e non affrontata in questo lavoro.

%  BIBLIOGRAFIA 

\printbibliography[heading=bibintoc,title={Bibliografia}]

\end{document}